\chapter{Translating expressions}
\label{chap:trexpr}

This chapter covers the layer that relates Lean's real kernel data structures, \texttt{Lean.Expr},
\texttt{Lean.LocalContext}, \texttt{Lean.Name} and \texttt{Lean.NameGenerator}, to the model syntax
\texttt{VExpr}/\texttt{VEnv} of Chapter \ref{chap:syntax}. Its job is to say precisely which real
expressions denote which model terms, and to prove that this correspondence survives every
operation the checker performs on a context: weakening by free or bound variables, instantiation of
a binder, abstraction of a free variable, universe substitution, environment growth and passage to a
definitionally equal context. The central relations are \texttt{TrExprS} (syntax-directed
translation, Definition \ref{ind:trexprs}), \texttt{TrExpr} (translation up to definitional
equality, Definition \ref{def:trexpr}) and, for structure projections, \texttt{TrProj}
(Definition \ref{def:trproj}).

Ten files are in scope.
\srcloc{Lean4Lean/Verify/Expr.lean}{1} (1219 lines) proves that the cached header bits of an
\texttt{Expr} agree with structural traversal, characterises application spines as lists, and
develops the lift/instantiate/abstract calculus.
\srcloc{Lean4Lean/Verify/Name.lean}{1} (90 lines) and
\srcloc{Lean4Lean/Verify/NameGenerator.lean}{1} (28 lines) supply the order properties of
\texttt{Name.quickCmp} and the freshness discipline.
\srcloc{Lean4Lean/Verify/LocalContext.lean}{1} (373 lines) models \texttt{mkBinding} as a fold and
translates a real local context.
\srcloc{Lean4Lean/Verify/VLCtx.lean}{1} (116 lines) defines the translation context.
\srcloc{Lean4Lean/Verify/Typing/Expr.lean}{1} (463 lines) holds the vocabulary: \texttt{Closed},
\texttt{FVarsIn}, \texttt{VLCtx.WF}, \texttt{TrProjCtor}/\texttt{TrProj}, \texttt{TrExprS}, the model
spellings of the literals and the primitive invariant \texttt{VEnv.HasPrimitives}.
\srcloc{Lean4Lean/Verify/Typing/Lemmas.lean}{1} (2357 lines) is essentially all of the metatheory of
the translation. The three consumers are
\srcloc{Lean4Lean/Verify/Typing/TrTerm.lean}{1} (204 lines, a bundled builder used by the primitive
proofs), \srcloc{Lean4Lean/Verify/Typing/ConditionallyTyped.lean}{1} (144 lines, the conditional
invariants the checker's monadic specifications are stated with) and
\srcloc{Lean4Lean/Verify/Typing/PrimSpec.lean}{1} (126 lines, environment extension for the
primitive invariant).

Nothing in this group is a thesis theorem: the thesis axiomatises the model, and this layer is the
bridge to the implementation, which the thesis does not discuss. Where the thesis does speak it
speaks about the model, and the master lemma families here are the translation-level shadows of its
results: \texttt{Closed}/\texttt{FVarsIn} and \texttt{TrExprS.closed}/\texttt{fvarsIn} of Regularity
(typesys.tex, \texttt{thm:reg}); the \texttt{FVLift'}/\texttt{BVLift} families and
\texttt{TrExprS.weakFV}/\texttt{weakBV} of Weakening (typesys.tex, \texttt{thm:weak}), including its
strengthening direction, which is exactly where the remaining holes sit; \texttt{InstN}/\texttt{InstLet}
and \texttt{TrExprS.instN}/\texttt{instN\_let} of Substitution (typesys.tex, \texttt{thm:subst});
and \texttt{TrExprS.uniq} of Unique typing (unique.tex, \texttt{thm:utype}).

The branches contributed 269 lines here. The \texttt{trproj} branch added 215: the whole of
\srcloc{Lean4Lean/Verify/Typing/Expr.lean}{69} to 135, which replaces master's one-line
\texttt{def TrProj ... := sorry} by \texttt{TrProjCtor}/\texttt{TrProj} (70 changed lines including
the two new parameters of the \texttt{TrExprS.proj} rule); 131 lines in
\srcloc{Lean4Lean/Verify/Typing/Lemmas.lean}{641}, namely the seven \texttt{TrProj} transport
lemmas, of which five are now proved and two remain \texttt{sorry}, the new \texttt{TrProj.mono},
\texttt{TrExprS.mkAppList\_inv}, and the one-line \texttt{proj} case of each master induction;
and 14 lines of spine lemmas in
\srcloc{Lean4Lean/Verify/Expr.lean}{745}. The \texttt{iota} branch added 54: a self-contained block
of lookup lemmas about the empty local context at
\srcloc{Lean4Lean/Verify/LocalContext.lean}{181} (42 lines), and ten lines of hypothesis
strengthening, eight signatures plus two call sites, from \texttt{env.Ordered} to
\texttt{env.OrderedStrong} in \srcloc{Lean4Lean/Verify/Typing/Lemmas.lean}{1856}, with two more
signatures in \srcloc{Lean4Lean/Verify/Typing/TrTerm.lean}{104}.

Assessment. The \texttt{trproj} work in this group is substantial and well judged: modelling a
projection as the recursor expansion of the thesis's \texttt{inv\_x}, carrying the expansion's
typing as a field rather than deriving it, and hiding the expansion data existentially, is the right
choice given that \texttt{VExpr} has no projection node, and it makes five of master's seven
\texttt{sorry} stubs go through with short, structured proofs while making \texttt{TrProj.wf}
unconditional. The two holes that remain, \texttt{TrProj.weak'\_inv} and \texttt{TrProj.uniq}, are
documented with precise statements of what is missing, and the first of them is blocked on a master
lemma. Two caveats should be carried forward: \texttt{TrProj} pins a pattern key rather than a
reduction rule, so it is weaker than it looks, and it has no general inhabitation theorem, since
\texttt{inferProj.WF\_struct} is \texttt{sorry}; and the \texttt{iota} hypothesis strengthening
quietly routes every \emph{use} of the literal lemmas, and hence the checker's literal handling,
through the admitted \texttt{VEnv.WF.patsStrong}.

Throughout, $\Gamma$ is a model typing context (a \texttt{List VExpr}), $\Delta$ a translation
context (\texttt{VLCtx}), $U$ a universe-parameter count, and we write $\Gamma \vdash e : A$ for
\texttt{VEnv.HasType} and $e_1 \equiv e_2$ for \texttt{VEnv.IsDefEqU} in the ambient environment.

\section{Real expressions}

\begin{lemma}[The cached \texttt{Expr} header bits agree with structural traversal]
\label{family:expr-flags}
\lean{Lean.Expr.hasFVar', Lean.Expr.containsFVar', Lean.Expr.hasFVar_eq, Lean.Expr.containsFVar_eq, Lean.Expr.hasExprMVar', Lean.Expr.hasExprMVar_eq, Lean.Expr.hasLevelMVar', Lean.Expr.hasLevelMVar_eq, Lean.Expr.hasLevelParam', Lean.Expr.hasLevelParam_eq, Lean.Expr.Data.looseBVarRange_le, Lean.Expr.mkData_looseBVarRange, Lean.Expr.mkAppData_looseBVarRange, Lean.Expr.hasFVar'_of_containsFVar'}
\leanok
Lean packs \texttt{hasFVar}, \texttt{hasExprMVar}, \texttt{hasLevelMVar}, \texttt{hasLevelParam} and
\texttt{looseBVarRange} into one machine word carried on every \texttt{Expr} node. Each is here
proved equal to the obvious structural recursion (the primed names), which is the form the rest of
the verification reasons with. Unglamorous but load-bearing: without it the verification could not
use the constant-time tests the real kernel uses.
\end{lemma}
\begin{proof}
\leanok
Structural inductions on the expression, with the bit-level facts about \texttt{Expr.Data}
(\texttt{mkData\_looseBVarRange}, \texttt{mkAppData\_looseBVarRange},
\texttt{Data.looseBVarRange\_le}) discharged by \texttt{bv\_decide}, whose native LRAT
certificates appear as extra axioms in the footprint.
\end{proof}
\axfoot{no sorryAx; bv\_decide LRAT certificates (Lean.Expr.mkData\_looseBVarRange.\_native.bv\_decide.ax\_1\_9 and siblings), plus Lean.Expr.mkData\_eq and Lean.Expr.mkAppData\_eq}

\begin{lemma}[Lawfulness of the \texttt{BEq}/\texttt{Ord} instances the kernel relies on]
\label{family:expr-beq}
\lean{Lean.Expr.eqv_refl, Lean.Expr.eqv_euc, Lean.Expr.eqv_sort, Lean.Expr.eqv_const, Lean.Expr.data_eq, Lean.Expr.instantiate1_eqv, Lean.Expr.instantiateList_eqv, Lean.Syntax.structEq_refl, Lean.Syntax.structEq_trans, Lean.Literal.beq_iff_eq, Lean.DefinitionSafety.compare, Lean.DefinitionSafety.le_antisymm, Lean.instLawfulBEqFVarId_lean4Lean, Lean.instLawfulBEqMVarId_lean4Lean, Lean.Expr.instEquivBEq_lean4Lean, Lean.Expr.instLawfulHashable_lean4Lean}
\leanok
\uses{family:expr-flags}
\texttt{FVarId}, \texttt{MVarId}, \texttt{LBool}, \texttt{DefinitionSafety}, \texttt{Substring.Raw},
\texttt{Syntax}, \texttt{DataValue}, \texttt{Literal} and \texttt{Expr} are given
\texttt{LawfulBEq}/\texttt{EquivBEq}/\texttt{TransOrd}/\texttt{LawfulHashable} instances, so that the
hash-map and comparison reasoning of the environment layer is sound. For \texttt{Expr} only
\texttt{EquivBEq} is available, since its \texttt{==} ignores binder names and binder information;
\texttt{LawfulHashable} then rests on \texttt{data\_eq}, which says equivalent expressions carry
equal header words.
\end{lemma}
\begin{proof}
\leanok
Simultaneous structural inductions over the two expressions; \texttt{eqv\_euc} (Euclidean
transitivity) is the shape the \texttt{EquivBEq} instance wants, and \texttt{data\_eq} then
follows by reading off the header word of each node.
\end{proof}

\begin{lemma}[Application spines as lists]
\label{family:expr-spine}
\lean{Lean.Expr.getAppArgsList, Lean.Expr.getAppArgsRevList, Lean.Expr.getAppArgsRevList_reverse, Lean.Expr.getAppNumArgs_eq, Lean.Expr.getAppArgs_eq, Lean.Expr.getAppRevArgs_eq, Lean.Expr.withApp_eq, Lean.Expr.withRevApp_eq, Lean.Expr.mkAppList, Lean.Expr.mkAppRevList, Lean.Expr.mkAppList_eq_foldl, Lean.Expr.mkAppRevList_eq_foldr, Lean.Expr.mkAppList_reverse, Lean.Expr.mkAppRevList_getAppArgsRevList, Lean.Expr.mkAppList_getAppArgsList, Lean.Expr.mkAppRange_eq, Lean.Expr.mkAppRange_eq_rev, Lean.Expr.mkAppRevRange_eq, Lean.Expr.mkAppRevRange_eq_rev}
\leanok
The array-based spine operations of \texttt{Lean.Expr} (\texttt{getAppArgs},
\texttt{getAppRevArgs}, \texttt{withApp}, \texttt{mkAppRange}) are characterised in terms of the
list-based \texttt{getAppArgsList}/\texttt{mkAppList}, which is the form the proofs use;
\texttt{mkAppList\_getAppArgsList} rebuilds an expression from its head and argument list.
\end{lemma}
\begin{proof}
\leanok
Structural inductions along the application spine, converting each array operation into its
list form and then commuting \texttt{reverse} through with
\texttt{getAppArgsRevList\_reverse}.
\end{proof}

\begin{theorem}[Reading back a spine built with \texttt{mkAppList}]
\label{thm:mkapplist-spine}
\lean{Lean.Expr.getAppFn_mkAppList, Lean.Expr.getAppArgsRevList_mkAppList, Lean.Expr.getAppArgsList_mkAppList}
\leanok
\contrib{trproj} \srcloc{Lean4Lean/Verify/Expr.lean}{745}
\thesisref{no counterpart; this is implementation-level bookkeeping}
For every list of arguments \texttt{es} and every expression \texttt{e}:
\texttt{(mkAppList e es).getAppFn = e.getAppFn}, and
\texttt{getAppArgsList (mkAppList e es) = getAppArgsList e ++ es}, with the reversed variant
\texttt{getAppArgsRevList (mkAppList e es) = es.reverse ++ getAppArgsRevList e} in between.
\end{theorem}
\begin{proof}
\uses{family:expr-spine}
\leanok
Induction on the argument list for the first two; the third rewrites through
\texttt{getAppArgsRevList\_reverse}. Three or four lines each.
\end{proof}

\begin{remark}
These three lemmas fill a real gap in master's spine API, which had only the
\texttt{mkAppList\_getAppArgsList} direction. Only one of the three has an external consumer:
\texttt{getAppArgsList\_mkAppList} is rewritten with at
\srcloc{Lean4Lean/Verify/TypeChecker/WHNF.lean}{50} in the projection and $\iota$ reduction proof.
\texttt{getAppFn\_mkAppList} is \texttt{@[simp]} and has no explicit call site, but it fires inside
the \texttt{simp} at \texttt{WHNF.lean:52}, where the argument list is not a literal.
\texttt{getAppArgsRevList\_mkAppList} has no consumer at all outside
\srcloc{Lean4Lean/Verify/Expr.lean}{749}, where it serves only to prove the forward form.
\end{remark}

\begin{lemma}[The substitution calculus of real Lean expressions]
\label{family:expr-subst}
\lean{Lean.Expr.liftLooseBVars_eq_self, Lean.Expr.liftLooseBVars_liftLooseBVars, Lean.Expr.liftLooseBVars_add, Lean.Expr.liftLooseBVars_comm, Lean.Expr.instantiate1'_looseBVarRange, Lean.Expr.instantiate1'_eq_self, Lean.Expr.instantiateList_eq_self, Lean.Expr.instantiate1'_liftLooseBVars, Lean.Expr.instantiate1'_instantiate1', Lean.Expr.instantiateList_eq_foldl, Lean.Expr.instantiateRevList_eq_foldr, Lean.Expr.instantiateList_instantiate1_comm, Lean.Expr.instantiateRev_eq_instantiateList, Lean.Expr.abstractList_eq_foldl, Lean.Expr.abstract1_comm, Lean.Expr.abstract1_abstractList, Lean.Expr.abstract1_hasLooseBVar, Lean.Expr.abstract1_eq_liftLooseBVars, Lean.Expr.abstract1_eq_self, Lean.Expr.abstractList_eq_self, Lean.Expr.lowerLooseBVars_eq_instantiate, Lean.Expr.abstract1_lower}
\leanok
Lifting, instantiation, abstraction and lowering of loose bound variables commute in the expected
ways, are the identity when the relevant range is empty, and can be written as folds over lists.
This is the source-side counterpart of the model's substitution calculus
(Chapter \ref{chap:syntax}), and the thesis's Substitution lemma (typesys.tex,
\texttt{thm:subst}) is its model-side analogue.
\end{lemma}
\begin{proof}
\uses{def:inst, def:liftn}
\leanok
Structural inductions on the expression with arithmetic side conditions on the indices; the
list-valued forms follow by a second induction on the substitution list. Of the family only
\texttt{instantiateRev\_eq\_instantiateList} reaches outside Lean's public API: it consumes the two
\texttt{Lean.Expr} implementation axioms \texttt{instantiateRev\_eq} and \texttt{instantiate\_eq}
of \srcloc{Lean4Lean/Verify/Axioms.lean}{1}.
\end{proof}

\begin{lemma}[Universe-parameter substitution on real expressions]
\label{family:expr-instlevel}
\lean{Lean.Expr.instantiateLevelParamsCore', Lean.Expr.instantiateLevelParamsCore_eq_self, Lean.Expr.instantiateLevelParamsCore_id, Lean.Expr.instantiateLevelParamsCore_eq, Lean.Expr.instantiateLevelParams_eq}
\leanok
\uses{family:expr-flags}
Lean's cached, short-circuiting \texttt{instantiateLevelParamsCore} agrees with the naive structural
substitution \texttt{instantiateLevelParamsCore'}, and is the identity on expressions with no level
parameter.
\end{lemma}
\begin{proof}
\uses{family:expr-flags}
\leanok
Induction on the expression, tracking the \texttt{hasLevelParam} bit: where the cached flag is
false the real implementation short-circuits and the naive recursion is the identity, which is
exactly \texttt{instantiateLevelParamsCore\_eq\_self}.
\end{proof}

\section{Names and the freshness discipline}

\begin{lemma}[Order properties of \texttt{Name.cmp} and \texttt{Name.quickCmp}]
\label{family:name-cmp}
\lean{Lean.Name.cmp_eq_swap, Lean.NameSet.contains_insert, Lean.NameSet.contains_empty, Lean.Name.instTransCmpCmp_lean4Lean, Lean.Name.instLawfulBEqCmpCmp_lean4Lean, Lean.Name.instTransCmpQuickCmp_lean4Lean, Lean.Name.instLawfulBEqCmpQuickCmp_lean4Lean}
\leanok
\texttt{Name.cmp} and the hash-accelerated \texttt{Name.quickCmp} are transitive and antisymmetric
comparators compatible with \texttt{==}, so that \texttt{NameMap}/\texttt{NameSet} (that is,
\texttt{Std.TreeSet}) reasoning is sound; \texttt{NameSet.contains\_insert} is the consequence
actually used downstream.
\end{lemma}
\begin{proof}
\leanok
\texttt{quickCmp} compares hashes before names, so transitivity is not obvious: the proof is a
simultaneous induction on the two names, reduced to \texttt{cmp\_eq\_swap} and an auxiliary
fact about the lexicographic \texttt{Ordering.then}.
\end{proof}

\begin{definition}[Freshness discipline of the name generator]
\label{def:ngen-reserves}
\lean{Lean.NameGenerator.Reserves, Lean.NameGenerator.next_reserves_self, Lean.NameGenerator.not_reserves_self, Lean.NameGenerator.LE, Lean.NameGenerator.LE.rfl, Lean.NameGenerator.LE.trans, Lean.NameGenerator.Reserves.mono, Lean.NameGenerator.LE.next}
\leanok
\texttt{ngen.Reserves fv} says that whenever \texttt{fv} is \texttt{ngen.namePrefix} numbered $i$,
then $i <$ \texttt{ngen.idx}: the generator has already handed that name out. Generators with the
same prefix are ordered by their index; \texttt{Reserves} is monotone along that order, the next
name is reserved by the successor generator and not by the current one. Small but structurally
important: this is the only thing that keeps the checker's fresh variables fresh, and
\texttt{ConditionallyTyped} (Definition \ref{def:conditionallytyped}) is stated in terms of it.
\end{definition}

\section{Real local contexts}

\begin{definition}[A list-based model of \texttt{LocalContext.mkBinding}]
\label{def:mkbindinglist}
\lean{Lean.LocalContext.mkBindingList1, Lean.LocalContext.mkBindingList, Lean.LocalContext.mkBinding_eq, Lean.LocalContext.mkBindingList1_abstract, Lean.LocalContext.mkBindingList_core_cons, Lean.LocalContext.mkBindingList_nil, Lean.LocalContext.mkBindingList_cons, Lean.LocalContext.mkBindingList_eq_fold, Lean.LocalContext.mkBindingList1_congr, Lean.LocalContext.mkBindingList_congr}
\leanok
Lean's \texttt{mkBinding} builds a lambda or pi telescope by abstracting an array of free variables,
using an internal loop and a cache. \texttt{mkBindingList} re-expresses it as a right fold of the
single-step \texttt{mkBindingList1}, which looks the variable up in the context and emits one
\texttt{lam}, \texttt{forallE} or \texttt{letE} binder (dropping an unused \texttt{let}).
\texttt{mkBinding\_eq} proves the two agree, and \texttt{mkBindingList\_congr} shows that only the
declarations actually looked up matter. The proofs are inductions on the variable list, unfolding
the real implementation's loop.
\end{definition}

\begin{definition}[Well-formed local contexts]
\label{ind:lctx-wf}
\lean{Lean.LocalContext.WF, Lean.LocalContext.WF.map_wf, Lean.LocalContext.WF.decls_wf, Lean.LocalContext.WF.map_toList, Lean.LocalContext.WF.find?_eq_find?_toList, Lean.LocalContext.WF.nodup, Lean.LocalContext.WF.mkLocalDecl, Lean.LocalContext.WF.mkLetDecl, Lean.LocalContext.mkLocalDecl_toList, Lean.LocalContext.mkLetDecl_toList}
\leanok
\texttt{LocalContext.WF} is the inductive predicate saying that a context is the empty one extended
by steps that insert a declaration under a variable absent from the map, with index equal to the
current array size. From it: the \texttt{PersistentArray} of declarations and the
\texttt{fvarIdToDecl} map agree as lists, \texttt{find?} equals lookup in \texttt{toList}, the
declared variables are pairwise distinct, and both \texttt{mkLocalDecl} and \texttt{mkLetDecl}
preserve well-formedness given freshness. The proofs are inductions on the construction, and they
consume the \texttt{PersistentArray}/\texttt{PersistentHashMap} axioms of
\srcloc{Lean4Lean/Verify/Axioms.lean}{1}.
\end{definition}

\begin{lemma}[Lookup in an explicitly built chain of local declarations]
\label{family:lctx-empty-decl}
\lean{Lean.LocalContext.WF.empty, Lean.LocalContext.toList_empty, Lean.LocalContext.find?_empty, Lean.LocalContext.empty_map_wf, Lean.LocalContext.map_wf_mkLocalDecl, Lean.LocalContext.find?_mkLocalDecl, Lean.LocalContext.instDecidableEqFVarId_lean4Lean}
\leanok
\contrib{iota} \srcloc{Lean4Lean/Verify/LocalContext.lean}{181}
\thesisref{no counterpart; support for the quotient axiom shapes}
The empty context is well formed, has empty \texttt{toList}, always-failing \texttt{find?} and a
well-formed \texttt{fvarIdToDecl}; declaring a variable preserves well-formedness of that map
\emph{whether or not} the variable is fresh; and
\texttt{(lctx.mkLocalDecl fv n ty bi k).find? x} is the new declaration if \texttt{x = fv} and
\texttt{lctx.find? x} otherwise, assuming only that the map is well formed. A
\texttt{DecidableEq FVarId} instance is declared alongside.
\end{lemma}
\begin{proof}
\uses{ind:lctx-wf}
\leanok
\texttt{find?\_mkLocalDecl} unfolds \texttt{find?} to a \texttt{PersistentHashMap.find?} on an
\texttt{insert}, applies \texttt{WF.find?\_insert} and splits on \texttt{x = fv};
\texttt{find?\_empty} goes through \texttt{WF.empty.find?\_eq\_find?\_toList} and
\texttt{toList\_empty}.
\end{proof}

\begin{remark}
The design point is real and is stated in the source: dropping the freshness side condition is what
makes a chain of six \texttt{mkLocalDecl} steps resolvable without discharging six freshness
obligations. Four of the seven declarations (\texttt{empty\_map\_wf},
\texttt{map\_wf\_mkLocalDecl}, \texttt{find?\_mkLocalDecl}, \texttt{find?\_empty}) are used,
heavily, in \srcloc{Lean4Lean/Verify/Environment/Quot.lean}{1}, where the \texttt{Quot} axiom
shapes are built by hand, and the \texttt{DecidableEq FVarId} instance is what makes the
\texttt{if x = fv} of \texttt{find?\_mkLocalDecl} reducible there; \texttt{WF.empty} and
\texttt{toList\_empty} serve only \texttt{find?\_empty} inside the block itself.
\end{remark}

\begin{theorem}[\texttt{mkForall} over declared distinct variables is a right fold]
\label{thm:mkforall-eq-fold}
\lean{Lean.LocalContext.mkForall_eq_fold}
\leanok
\contrib{iota} \srcloc{Lean4Lean/Verify/LocalContext.lean}{217}
\thesisref{no counterpart}
For a list \texttt{xs} of pairwise distinct free variables, all declared in \texttt{lctx}, and any
body \texttt{b}, \texttt{lctx.mkForall (xs.map .fvar) b} equals the right fold over \texttt{xs} of
\texttt{fun a e => mkBindingList1 false lctx [] a (e.abstract1 a)}, that is, one binder abstracted
at a time.
\end{theorem}
\begin{proof}
\uses{def:mkbindinglist}
\leanok
One line: \texttt{mkBinding\_eq} reduces \texttt{mkForall} to \texttt{mkBindingList} and master's
\texttt{mkBindingList\_eq\_fold} to the fold; the hypotheses are exactly the latter's. Thin, but it
is rewritten with eight times in \srcloc{Lean4Lean/Verify/Environment/Quot.lean}{163} to
\texttt{:238}, to compute the shapes of the \texttt{Quot} axioms.
\end{proof}

\section{Translation contexts}

\begin{definition}[Translated local declarations]
\label{ind:vt-vlocaldecl}
\lean{Lean4Lean.VLocalDecl, Lean4Lean.VLocalDecl.depth, Lean4Lean.VLocalDecl.value, Lean4Lean.VLocalDecl.type, Lean4Lean.VLocalDecl.type', Lean4Lean.VLocalDecl.lift', Lean4Lean.VLocalDecl.liftN, Lean4Lean.VLocalDecl.inst, Lean4Lean.VLocalDecl.instL}
\leanok
\uses{ind:vexpr}
A \texttt{VLocalDecl} is either \texttt{vlam type}, a real binder, or \texttt{vlet type value}, a
let-binding. Its \texttt{depth} is $1$ for \texttt{vlam} and $0$ for \texttt{vlet}; its
\texttt{value} is \texttt{bvar 0} for \texttt{vlam} and the definition for \texttt{vlet}; its
\texttt{type} lifts in the \texttt{vlam} case. Lifting, instantiation and level instantiation act
componentwise. The \texttt{depth} field is the device that lets a \texttt{let} occupy a slot in the
source context while vanishing from the model context.
\end{definition}

\begin{definition}[The translation context and variable lookup]
\label{def:vt-vlctx}
\lean{Lean4Lean.VLCtx, Lean4Lean.VLCtx.bvars, Lean4Lean.VLCtx.NoBV, Lean4Lean.VLCtx.next, Lean4Lean.VLCtx.find?, Lean4Lean.VLCtx.liftVar, Lean4Lean.VLCtx.varToExpr, Lean4Lean.VLCtx.vlamName, Lean4Lean.VLCtx.fvars, Lean4Lean.VLCtx.toCtx, Lean4Lean.VLCtx.instL}
\leanok
\uses{ind:vt-vlocaldecl}
A \texttt{VLCtx} is a list of pairs of an optional name (an \texttt{FVarId} with the list of free
variables its declaration depends on) and a \texttt{VLocalDecl}.
\texttt{find? $\Delta$ v} resolves a source variable \texttt{v}, either a de Bruijn index or a free
variable, to a pair (model term, model type), lifting by the accumulated \texttt{depth} as it walks
outwards; \texttt{toCtx} is the model typing context, which drops \texttt{vlet} entries;
\texttt{fvars} lists the named entries and \texttt{bvars} counts the anonymous ones. This is the
pivot of the whole \texttt{Verify} layer: it is simultaneously a source-side scope and a model-side
context, and the mismatch between the two is what the lift and instantiation relations below manage.
\end{definition}

\begin{lemma}[Basic lookup lemmas for translation contexts]
\label{family:vt-vlctx-basic}
\lean{Lean4Lean.VLCtx.fvars_nil, Lean4Lean.VLCtx.fvars_cons_none, Lean4Lean.VLCtx.fvars_cons_some, Lean4Lean.VLCtx.find?_eq_some, Lean4Lean.VLCtx.vlamName_mem_fvars}
\leanok
\texttt{fvars} computes on the two \texttt{cons} shapes; a free variable has a successful
\texttt{find?} if and only if it is in \texttt{fvars}; and a \texttt{vlamName} hit at index $i$ is a
member of \texttt{fvars}. Structural inductions on the context.
\end{lemma}
\begin{proof}
\uses{def:vt-vlctx}
\leanok
Induction on $\Delta$, unfolding \texttt{find?} and \texttt{next} at each step and splitting on
whether the head entry's name is the one sought.
\end{proof}

\section{Closedness and free variables}

\begin{definition}[Closedness of a real Lean expression]
\label{def:vt-closed}
\lean{Lean4Lean.Closed, Lean.Expr.Closed}
\leanok
\texttt{Closed e k} says that the loose de Bruijn indices of \texttt{e} are all below \texttt{k},
incrementing \texttt{k} at each binder, \emph{and} that \texttt{e} contains no expression
metavariable: the \texttt{.mvar} case is \texttt{False}, not \texttt{True}. The asymmetry is
deliberate and makes \texttt{Closed} strictly stronger than a bound on the loose-variable range;
\texttt{Closed.of\_looseBVarRange} is where the difference has to be paid for. It is the syntactic
side of the thesis's Regularity lemma (typesys.tex, \texttt{thm:reg}), clause
$FV(e) \subseteq \Gamma$.
\end{definition}

\begin{definition}[Specification-level list of free variables]
\label{def:vt-fvarslist}
\lean{Lean.Expr.fvarsList}
\leanok
\texttt{e.fvarsList} collects every \texttt{FVarId} occurring in \texttt{e}, in traversal order and
with duplicates. Its docstring marks it as a specification-only device: it is quadratic and is never
run.
\end{definition}

\begin{definition}[Free variables of an expression lie in a predicate]
\label{def:vt-fvarsin}
\lean{Lean4Lean.FVarsIn, Lean.Expr.FVarsIn}
\leanok
\uses{def:vt-fvarslist}
\texttt{FVarsIn P e} says that every free variable of \texttt{e} satisfies \texttt{P}, that
\texttt{e} has no expression metavariable, and that it has no \emph{level} metavariable either: the
\texttt{sort} and \texttt{const} cases require \texttt{hasMVar' = false} on the levels they carry.
It is the invariant the checker carries about which variables a term may mention. Bundling three
different conditions into one predicate is convenient at the call sites but makes \texttt{FVarsIn}
do double duty; \texttt{fvarsIn\_iff} splits it again.
\end{definition}

\begin{lemma}[Calculus of closedness]
\label{family:vt-closed-lemmas}
\lean{Lean4Lean.Closed.mono, Lean4Lean.Closed.natLitToConstructor, Lean4Lean.Closed.strLitToConstructor, Lean4Lean.Closed.toConstructor, Lean4Lean.Closed.litType, Lean4Lean.Closed.abstract1, Lean4Lean.Closed.getAppFn, Lean4Lean.Closed.getAppArgsRevList, Lean4Lean.Closed.getAppArgsList, Lean4Lean.Closed.looseBVarRange_le, Lean4Lean.Closed.looseBVarRange_zero, Lean4Lean.Closed.of_looseBVarRange, Lean4Lean.Closed.of_looseBVarRange_zero}
\leanok
\texttt{Closed} is monotone in the bound, holds of the constructor forms of literals, is preserved
by abstraction and by taking the head or the arguments of an application spine, and is
interconvertible with the cached \texttt{looseBVarRange} provided the term has no expression
metavariable.
\end{lemma}
\begin{proof}
\uses{def:vt-closed, family:expr-flags}
\leanok
Structural inductions on the expression. The direction from \texttt{looseBVarRange} back to
\texttt{Closed} is the one that needs the metavariable hypothesis, as its docstring notes.
\end{proof}

\begin{lemma}[Calculus of the free-variable predicate]
\label{family:vt-fvarsin-lemmas}
\lean{Lean4Lean.fvarsIn_iff, Lean4Lean.fvarsIn_iff_hasMVar, Lean4Lean.fvarsList_eq_nil, Lean4Lean.FVarsIn.mp, Lean4Lean.FVarsIn.mono, Lean4Lean.FVarsIn.fvars_cons, Lean4Lean.FVarsIn.default, Lean4Lean.FVarsIn.getRevArg!, Lean4Lean.FVarsIn.getArg!, Lean4Lean.FVarsIn.abstract_instantiate1, Lean4Lean.FVarsIn.abstract_eq_self, Lean4Lean.FVarsIn.of_abstract1, Lean4Lean.FVarsIn.of_abstractList, Lean4Lean.FVarsIn.liftLooseBVars, Lean4Lean.FVarsIn.instantiate1, Lean4Lean.FVarsIn.instantiateList, Lean4Lean.FVarsIn.abstract1, Lean4Lean.FVarsIn.appRevList, Lean4Lean.FVarsIn.abstract1_erase, Lean4Lean.FVarsIn.of_hasFVar, Lean4Lean.FVarsIn.of_containsFVar', Lean4Lean.FVarsIn.natLitToConstructor, Lean4Lean.FVarsIn.strLitToConstructor, Lean4Lean.FVarsIn.toConstructor, Lean4Lean.FVarsIn.litType}
\leanok
\texttt{FVarsIn} is characterised by \texttt{fvarsList} together with the two metavariable flags, is
monotone and closed under pointwise implication, is preserved by lifting, instantiation, abstraction
and spine building, and can be established from the cached \texttt{hasFVar}, \texttt{hasExprMVar}
and \texttt{hasLevelMVar} bits.
\end{lemma}
\begin{proof}
\uses{def:vt-fvarsin, def:vt-fvarslist, family:expr-flags, family:expr-subst}
\leanok
Structural inductions on the expression, with \texttt{simp} and \texttt{grind} closing the routine
cases; the flag-based lemmas go through the header-bit characterisations.
\end{proof}

\section{Well-formed translation contexts}

\begin{definition}[Well-formedness of a single translated local declaration]
\label{def:vt-vlocaldecl-wf}
\lean{Lean4Lean.VLocalDecl.WF}
\leanok
\uses{ind:vt-vlocaldecl, def:hastype-istype}
A \texttt{vlam type} entry is well formed over $\Gamma$ when \texttt{type} is a type there, and a
\texttt{vlet type value} entry when $\Gamma \vdash \texttt{value} : \texttt{type}$.
\end{definition}

\begin{definition}[Well-formedness of a translation context]
\label{def:vt-vlctx-wf}
\lean{Lean4Lean.VLCtx.FVWF, Lean4Lean.VLCtx.WF, Lean4Lean.VLCtx.WF.fvwf}
\leanok
\uses{def:vt-vlocaldecl-wf, def:vt-vlctx}
\texttt{VLCtx.FVWF} requires of every named entry that its \texttt{FVarId} is absent from the tail's
\texttt{fvars} and that its recorded dependencies are contained in them.
\texttt{VLCtx.WF env U} requires the same plus well-formedness of each declaration over the tail's
\texttt{toCtx}; \texttt{WF.fvwf} forgets the typing half. The name condition alone is what most
scope lemmas need, which is why the two are separated.
\end{definition}

\begin{lemma}[Structural lemmas for translated local declarations]
\label{family:vt-vlocaldecl-wf-lemmas}
\lean{Lean4Lean.VLocalDecl.lift'_consN_skipN, Lean4Lean.VLocalDecl.WF.hasType, Lean4Lean.VLocalDecl.WF.weakN, Lean4Lean.VLocalDecl.WF.instN, Lean4Lean.VLocalDecl.WF.instL, Lean4Lean.VLocalDecl.is_liftN, Lean4Lean.VLocalDecl.IsDefEq, Lean4Lean.VLocalDecl.lift'_depth, Lean4Lean.VLocalDecl.lift'_comp, Lean4Lean.VLocalDecl.weak'_iff, Lean4Lean.VLocalDecl.weakN_iff}
\leanok
A well-formed declaration types its \texttt{value} at its \texttt{type} in the extended context;
well-formedness is preserved by weakening, instantiation and level instantiation, and, given unique
typing, reflected by weakening (\texttt{weak'\_iff}, \texttt{weakN\_iff}).
\texttt{VLocalDecl.IsDefEq} is the pointwise definitional-equality relation on declarations.
\end{lemma}
\begin{proof}
\uses{def:vt-vlocaldecl-wf, thm:isdefeq-weakn, thm:isdefeq-instn, thm:isdefeq-instl, thm:isdefequ-weakn-iff}
\leanok
Case split on \texttt{vlam}/\texttt{vlet} and appeal to the corresponding model lemma.
\end{proof}
\axfoot{sorryAx in VLocalDecl.weak'\_iff and VLocalDecl.weakN\_iff only, inherited from VEnv.IsDefEqU.weakN\_iff (Theory/Typing/UniqueTyping.lean:172, master)}

\begin{lemma}[Well-formed translation contexts]
\label{family:vt-vlctx-wf-lemmas}
\lean{Lean4Lean.VLCtx.WF.find?_wf, Lean4Lean.VLCtx.WF.toCtx, Lean4Lean.VLCtx.WF.fvars_nodup, Lean4Lean.VLCtx.liftVar_zero, Lean4Lean.VLCtx.instL_eq_map, Lean4Lean.VLCtx.instL_toCtx, Lean4Lean.VLCtx.WF.instL, Lean4Lean.VLCtx.find?_instL, Lean4Lean.VLCtx.SortList}
\leanok
In a well-formed context every successful \texttt{find?} yields a typing derivation for the pair it
returns, \texttt{toCtx} is a well-formed model context, the free variables are pairwise distinct,
and both well-formedness and lookup commute with level instantiation. \texttt{SortList} records a
list of levels sorting the types of the outermost named binders.
\end{lemma}
\begin{proof}
\uses{def:vt-vlctx-wf}
\leanok
Inductions on the context; \texttt{find?\_wf} lifts the declaration's typing by the accumulated
depth at each step.
\end{proof}

\begin{definition}[Definitional equality of translation contexts]
\label{ind:vt-vlctx-isdefeq}
\lean{Lean4Lean.VLCtx.IsDefEq, Lean4Lean.VLocalDecl.IsDefEq.refl, Lean4Lean.VLCtx.IsDefEq.refl, Lean4Lean.VLCtx.IsDefEq.defeqCtx, Lean4Lean.VLCtx.IsDefEq.fvars, Lean4Lean.VLocalDecl.IsDefEq.wf, Lean4Lean.VLCtx.IsDefEq.wf, Lean4Lean.VLocalDecl.IsDefEq.symm, Lean4Lean.VLocalDecl.IsDefEq.defeqDFC, Lean4Lean.VLCtx.IsDefEq.symm, Lean4Lean.VLCtx.IsDefEq.find?_uniq, Lean4Lean.VLCtx.IsDefEq.find?_defeqDFC}
\leanok
\uses{def:vt-vlctx-wf, fam:defeqdfc}
Two translation contexts are definitionally equal when they have the same shape, the same names, and
pointwise definitionally equal declarations. Such a pair induces an \texttt{IsDefEqCtx} on the model
contexts, has the same \texttt{fvars}, is symmetric, is well formed on both sides, and makes
\texttt{find?} return definitionally equal answers on the two sides
(\texttt{find?\_uniq}, \texttt{find?\_defeqDFC}). Inductions on the relation and on \texttt{find?}.
\end{definition}

\section{The projection translation}

\begin{definition}[Projection as a recursor expansion, \texttt{TrProjCtor}]
\label{struct:trprojctor}
\lean{Lean4Lean.TrProjCtor}
\leanok
\uses{def:proj-fn, def:proj-motive-body, def:proj-inst-pis, def:pi-telescope, def:ctorheaded, def:proj-binder-arity, inductive:simple-pattern, def:hastype-istype, def:mkapps}
\contrib{trproj} \srcloc{Lean4Lean/Verify/Typing/Expr.lean}{90}
\thesisref{typesys.tex, undecidability of definitional equality, the \texttt{inv\_x} construction; Wtypes.tex for the dependent motive of $\pi_2$; axioms.tex 2.6.3--2.6.4 for the recursor and its $\iota$ rule}
\texttt{TrProjCtor env U $\Gamma$ S i e e' ctorName usS uss params np fieldTys} relates a model term
\texttt{e} of type \texttt{S usS params} to the translation \texttt{e'} of the source projection
\texttt{.proj S i e}. Since \texttt{VExpr} has no projection node, \texttt{e'} is required to be
\texttt{.app (VExpr.projFn S usS uss params fieldTys i) e}. Its eight fields require, in
order:
some reduct is registered in \texttt{env.pats} for the $\iota$ key
\texttt{SimplePattern.iota (mkRecName S) (np+1+1+0) ctorName (np + fieldTys.length)};
\texttt{params.length = np}; the constructor \texttt{ctorName} exists, its recorded type is
\texttt{CtorHeaded}, and \texttt{fieldTys} is the $\Pi$-binder telescope of that type instantiated at
\texttt{usS} and \texttt{params}; \texttt{i < fieldTys.length}; the recursor constant exists and its binder number
\texttt{np+1} has arity \texttt{fieldTys.length}; $\Gamma \vdash e : $ \texttt{S usS params}; the
expansion has type \texttt{forallE (S usS params) (projMotiveBody S usS uss params fieldTys i)};
and \texttt{e'} is that expansion applied to \texttt{e}.
\end{definition}

\begin{remark}
Three points about the interface. (1) The \texttt{pat} field asserts only that \emph{some} reduct is
registered under that key: the reduct itself is not pinned, so \texttt{TrProjCtor} alone does not
determine reduction behaviour, and the key, being a sum, does not by itself exclude reading a
two-minor recursor as \texttt{(np+1)+1+1+0}. The docstring says so and defers to
\texttt{TrEnv.proj\_defeq} (Theorem \ref{thm:tr-env-proj-defeq}) for the kernel-side facts.
(2) \texttt{uss} is a total function \texttt{Nat} to \texttt{List VLevel}, so the relation quantifies
over a whole level assignment of which only the first \texttt{i+1} values matter; this is what makes
\texttt{TrProj.uniq} hard. (3) The scope paragraph, which restricts the relation to single
constructor, non-recursive, non-indexed, non-mutual structures and declares reflexive, indexed and
nested structures a completeness rather than a soundness boundary, makes precise claims about the
kernel's \texttt{inferProj} that are prose only; nothing in the repository formalizes them.
Documentation density is otherwise unusually high, with a per-field docstring naming the kernel fact
each field stands in for.
\end{remark}

\begin{definition}[The projection-translation relation \texttt{TrProj}]
\label{def:trproj}
\lean{Lean4Lean.TrProj}
\leanok
\uses{struct:trprojctor}
\contrib{trproj} \srcloc{Lean4Lean/Verify/Typing/Expr.lean}{133}
\thesisref{as for Definition \ref{struct:trprojctor}}
\texttt{TrProj env U $\Gamma$ S i e e'} existentially quantifies the constructor name, the
structure's level list \texttt{usS}, the per-field level assignment \texttt{uss} (a function, not a
list), the parameters, the parameter count and the field telescope of
Definition \ref{struct:trprojctor}: there is \emph{some} single-constructor reading of \texttt{S}
under which \texttt{e'} is the recursor expansion of the \texttt{i}-th projection of \texttt{e}. On
master this was \texttt{def TrProj : ... := sorry}; the branch keeps the stub's argument order and
adds the two parameters \texttt{env} and \texttt{U} that the typing fields need, which is what forces
the new \texttt{TrProj.mono} (Theorem \ref{thm:trproj-mono}) and the extra arguments in the
\texttt{TrExprS.proj} rule.
\end{definition}

\begin{remark}
Two of the six existentials are pinned by other fields: \texttt{np} is \texttt{params.length} by
\texttt{params\_length}, and \texttt{fieldTys} is determined by \texttt{ctorName}, \texttt{usS} and
\texttt{params} by \texttt{ctor}. Replacing them by definitions would shrink the relation and remove
an obligation from each of the three proofs that rebuild the record
(\texttt{weak'}, \texttt{instN}, \texttt{instL}), where they reappear as
\texttt{params\_length := by simpa} and \texttt{field\_lt := by rw [hlen]};
\texttt{defeqDFC} and \texttt{mono}, which reuse the old record with \texttt{\{H with \ldots\}},
pay nothing.
\end{remark}

\section{The translation relation}

\begin{definition}[Environment support needed by a literal]
\label{def:vt-containslits}
\lean{Lean4Lean.VEnv.ContainsLits}
\leanok
A \texttt{Nat} literal needs \texttt{Nat} in the environment; a string literal needs
\texttt{Char.ofNat} and \texttt{String.ofList}. This is the side condition of the \texttt{lit} rule
below, and the reason literals do not silently translate in an environment that cannot express them.
\end{definition}

\begin{definition}[Strict translation of a Lean expression, \texttt{TrExprS}]
\label{ind:trexprs}
\lean{Lean4Lean.TrExprS}
\leanok
\uses{def:trproj, def:vt-vlctx, def:vt-containslits, def:vlevel-oflevel, def:hastype-istype}
\contrib{mixed} \srcloc{Lean4Lean/Verify/Typing/Expr.lean}{142}
\thesisref{no counterpart; the thesis does not model the implementation's expressions}
\texttt{TrExprS env Us $\Delta$ e e'} is the syntax-directed relation translating a real
\texttt{Lean.Expr} into a \texttt{VExpr} over $\Delta$: variables of either kind go through
\texttt{VLCtx.find?}, sorts and constant levels through \texttt{VLevel.ofLevel Us},
\texttt{mdata} is erased, \texttt{letE} is substituted away by translating the body in the context
extended with a \texttt{vlet} entry, a literal is translated through its constructor form, and
\texttt{.proj s i e} is translated by \texttt{TrProj} applied to the translation of \texttt{e}. The
\texttt{app}, \texttt{lam}, \texttt{forallE} and \texttt{letE} rules carry the model typing side
conditions needed to rebuild typing derivations.
\end{definition}

\begin{remark}
Master's inductive; the \texttt{trproj} branch changed the \texttt{proj} rule's premise from
\texttt{TrProj $\Delta$.toCtx s i e' e''} to
\texttt{TrProj env Us.length $\Delta$.toCtx s i e' e''}. The change is forced by \texttt{TrProj}
now carrying typing premises, and it means that in the projection case \texttt{TrExprS} is no longer
environment-independent, which is what creates the obligation \texttt{TrProj.mono}.
\end{remark}

\begin{definition}[Translation up to definitional equality, \texttt{TrExpr}]
\label{def:trexpr}
\lean{Lean4Lean.TrExpr}
\leanok
\uses{ind:trexprs, def:isdefeq}
\texttt{TrExpr env Us $\Delta$ e e'} holds when there is an \texttt{e2} with
\texttt{TrExprS env Us $\Delta$ e e2} and $\texttt{e2} \equiv \texttt{e'}$ over
$\Delta$\texttt{.toCtx}. This is the relation the checker's specifications are stated with, since a
checker is entitled to return any definitionally equal answer.
\end{definition}

\begin{definition}[Translating a real local context to a \texttt{VLCtx}]
\label{ind:trlctx}
\lean{Lean4Lean.TrLocalDecl, Lean4Lean.TrLocalDecl.wf, Lean.LocalDecl.deps, Lean4Lean.TrLocalDecl.deps_wf, Lean4Lean.TrLCtx', Lean4Lean.TrLCtx, Lean4Lean.TrLCtx.nil, Lean4Lean.TrLCtx'.noBV, Lean4Lean.TrLCtx'.fvars_eq, Lean4Lean.TrLCtx.fvars_eq, Lean4Lean.TrLCtx'.find?_eq_some, Lean4Lean.TrLCtx.find?_eq_some, Lean4Lean.TrLCtx.find?_eq_none, Lean4Lean.TrLCtx.contains, Lean4Lean.TrLCtx'.wf, Lean4Lean.TrLCtx.wf, Lean.LocalDecl.value', Lean4Lean.TrLCtx'.find?_of_mem, Lean4Lean.TrLCtx.find?_of_mem, Lean4Lean.TrLCtx.mkLocalDecl, Lean4Lean.TrLCtx.mkLetDecl, Lean4Lean.TrLCtx'.forall₂}
\leanok
\uses{ind:lctx-wf, def:vt-vlctx-wf, ind:trexprs}
\texttt{TrLCtx env Us lctx $\Delta$} says that \texttt{lctx} is well formed and that its declaration
list translates entry by entry to $\Delta$: a \texttt{cdecl} to a \texttt{vlam}, an \texttt{ldecl}
to a \texttt{vlet}, each named entry recording the declaration's \texttt{deps}, the free variables of
its type (and value). Consequences: the \texttt{fvars} agree as lists, lookups correspond, $\Delta$
has no bound variables, $\Delta$ is well formed, and \texttt{mkLocalDecl}/\texttt{mkLetDecl} extend
the two sides in step. This is the bridge between the checker's real \texttt{LocalContext} and the
model context. Inductions on the declaration list.
\end{definition}

\begin{theorem}[A translatable expression is scoped by its context]
\label{thm:trexprs-closed-fvars}
\lean{Lean4Lean.TrExprS.closed, Lean4Lean.TrExprS.fvarsIn, Lean4Lean.TrExprS.fvarsList, Lean4Lean.TrExpr.closed, Lean4Lean.TrExpr.fvarsIn, Lean4Lean.TrExpr.fvarsList, Lean4Lean.ofLevel_hasMVar}
\leanok
\uses{ind:trexprs, def:vt-closed, def:vt-fvarsin}
If \texttt{TrExprS env Us $\Delta$ e e'} then \texttt{e} is \texttt{Closed} under
$\Delta$\texttt{.bvars}, and \texttt{FVarsIn} the membership predicate of $\Delta$\texttt{.fvars},
so in particular \texttt{e} has no metavariable of either kind. The same holds for \texttt{TrExpr}.
This is the translation-level form of the thesis's Regularity lemma
(typesys.tex, \texttt{thm:reg}).
\end{theorem}
\begin{proof}
\uses{family:vt-vlctx-basic, family:vt-fvarsin-lemmas}
\leanok
Induction on the translation. The bound-variable case is an inner induction on the context showing
that a successful \texttt{find?} at index $i$ forces $i < \Delta$\texttt{.bvars}; the level cases
use \texttt{ofLevel\_hasMVar}.
\end{proof}

\section{Weakening, instantiation and abstraction of translation contexts}

\begin{definition}[Weakening a translation context by free variables]
\label{ind:vt-fvlift}
\lean{Lean4Lean.VLCtx.FVLift', Lean4Lean.VLCtx.FVLift'.toCtx, Lean4Lean.VLCtx.FVLift'.comp, Lean4Lean.VLCtx.FVLift'.from_nil, Lean4Lean.VLCtx.FVLift'.fvars_sublist, Lean4Lean.VLCtx.FVLift'.bvars_eq, Lean4Lean.VLCtx.FVLift'.wf, Lean4Lean.VLCtx.FVLift'.find?, Lean4Lean.VLCtx.FVLift, Lean4Lean.VLCtx.FVLift.toFVLift', Lean4Lean.VLCtx.FVLift.toCtx, Lean4Lean.VLCtx.FVLift.from_nil, Lean4Lean.VLCtx.FVLift.wf, Lean4Lean.VLCtx.FVLift.fvars_suffix, Lean4Lean.VLCtx.FVLift.find?}
\leanok
\uses{def:vt-vlctx-wf, fam:ctx-lift}
\texttt{FVLift' $\Delta$ $\Delta'$ dk n k} records that $\Delta'$ is obtained from $\Delta$ by
inserting named declarations (\texttt{skip\_fvar}), keeping named ones (\texttt{cons\_fvar}) and
keeping anonymous ones (\texttt{cons\_bvar}), while tracking the induced lift \texttt{n} of the model
context; \texttt{FVLift} is the simpler variant in which no anonymous entry is crossed. The
accompanying lemmas transport the relation to a \texttt{Ctx.Lift'} on \texttt{toCtx}, compose two
lifts, show \texttt{fvars} is a sublist and \texttt{bvars} unchanged, push well-formedness downwards,
and send a successful \texttt{find?} in $\Delta$ to the lifted answer in $\Delta'$. This is the
translation-level form of Weakening (typesys.tex, \texttt{thm:weak}). Inductions on the relation
with a \texttt{vlam}/\texttt{vlet} case split at each step.
\end{definition}

\begin{definition}[Weakening a translation context by bound variables]
\label{ind:vt-bvlift}
\lean{Lean4Lean.VLCtx.BVLift, Lean4Lean.VLCtx.BVLift.toCtx, Lean4Lean.VLCtx.BVLift.wf, Lean4Lean.VLCtx.BVLift.fvars_eq, Lean4Lean.VLCtx.BVLift.find?}
\leanok
\uses{ind:vt-fvlift}
\texttt{BVLift} is the analogue that inserts anonymous (bound-variable) entries: the source term must
then be lifted with \texttt{liftLooseBVars'} while the model term is lifted with \texttt{liftN}, and
\texttt{find?} commutes with \texttt{liftVar}. The \texttt{fvars} are unchanged. Induction on the
relation.
\end{definition}

\begin{definition}[Instantiating a binder of a translation context]
\label{ind:vt-instn}
\lean{Lean4Lean.VLCtx.InstN, Lean4Lean.VLCtx.InstN.toCtx, Lean4Lean.VLCtx.InstN.wf, Lean4Lean.VLCtx.InstN.fvars_eq, Lean4Lean.VLCtx.InstLet, Lean4Lean.VLCtx.InstLet.toCtx, Lean4Lean.VLCtx.InstLet.wf, Lean4Lean.VLCtx.InstLet.fvars_eq}
\leanok
\uses{ind:vt-fvlift, fam:ctx-lift}
\texttt{InstN $\Delta_0$ e$_0$ A$_0$ dk k $\Delta_1$ $\Delta$} says that $\Delta$ arises from
$\Delta_1$ by substituting \texttt{e}$_0$ for the anonymous \texttt{vlam} binder at depth \texttt{dk},
with \texttt{k} the corresponding model position; \texttt{InstLet} is the analogue for a
\texttt{vlet} binder, which leaves \texttt{toCtx} unchanged. Both preserve well-formedness (given a
typing of \texttt{e}$_0$) and the free-variable list, and both map to \texttt{Ctx.InstN} on
\texttt{toCtx}. This is the translation-level form of Substitution
(typesys.tex, \texttt{thm:subst}).
\end{definition}

\begin{definition}[Turning a free variable of a translation context back into a binder]
\label{ind:vt-abstract}
\lean{Lean4Lean.VLCtx.Abstract, Lean4Lean.VLCtx.Abstract.toCtx, Lean4Lean.VLCtx.Abstract.wf, Lean4Lean.VLCtx.Abstract.fvars_eq, Lean4Lean.VLCtx.Abstract.find?_self, Lean4Lean.VLCtx.Abstract.find?}
\leanok
\uses{ind:vt-instn}
\texttt{Abstract $\Delta_0$ v$_0$ d$_0$ dk k $\Delta_1$ $\Delta$} relates a context whose relevant
entry is the named declaration \texttt{v}$_0$ to the same context with that entry made anonymous.
The model context is unchanged, the free variables lose exactly \texttt{v}$_0$, \texttt{find?}
agrees away from \texttt{v}$_0$, and \texttt{v}$_0$ itself now resolves to the new bound variable.
Induction on the relation.
\end{definition}

\section{Transport of the projection translation}

The seven lemmas of this section were all \texttt{sorry} on master; five are proved here, two remain
open, and \texttt{TrProj.mono} is a new obligation created by the design.

\begin{theorem}[Weakening of the projection translation]
\label{thm:trproj-weak}
\lean{Lean4Lean.TrProj.weak', Lean4Lean.TrProj.weakN}
\leanok
\uses{def:trproj, fam:ctx-lift}
\contrib{trproj} \srcloc{Lean4Lean/Verify/Typing/Lemmas.lean}{641}
\thesisref{Weakening (typesys.tex, \texttt{thm:weak}), at the level of the translation}
If \texttt{env} is \texttt{Ordered}, \texttt{W : Ctx.Lift' n $\Gamma$ $\Gamma'$} and
\texttt{TrProj env U $\Gamma$ s i e e'}, then
\texttt{TrProj env U $\Gamma'$ s i (e.lift' n) (e'.lift' n)}; \texttt{weakN} is the
\texttt{Ctx.LiftN} form, obtained by specialising the general lift. The \texttt{Ordered}
hypothesis is new relative to master's stub, which took none.
\end{theorem}
\begin{proof}
\uses{fam:proj-fn-subst, fam:proj-motivebody-subst, fam:pibinders-stability, fam:proj-instpis-subst, thm:isdefeq-weakn, thm:isdefeq-closedn}
\leanok
One \texttt{refine} over a record literal: \texttt{params} is lifted and the field telescope lifted
pointwise by \texttt{piBinders\_lift'}. The key, the field index and the minor arity survive
because they mention \texttt{fieldTys} only through its length, so \texttt{List.length\_mapIdx}
discharges them, and the parameter count by \texttt{List.length\_map}; the two typings are weakened
by \texttt{HasType.weak'} and realigned by \texttt{projFn\_lift'} and
\texttt{projMotiveBody\_lift'}; and the constructor-instantiation equation is transported by
\texttt{instPis\_lift'} together with closedness of the recorded constant type
(\texttt{Ordered.closedC}); the \texttt{Ordered} hypothesis is consumed there and by the two
weakenings, \texttt{HasType.weak'} itself requiring it.
\end{proof}

\begin{theorem}[Strengthening of the projection translation (open)]
\label{thm:trproj-weak-inv}
\lean{Lean4Lean.TrProj.weak'_inv}
\leanok
\uses{def:trproj, fam:ctx-lift}
\contrib{trproj} \srcloc{Lean4Lean/Verify/Typing/Lemmas.lean}{745}
\thesisref{Weakening (typesys.tex, \texttt{thm:weak}), item 3, the strengthening direction}
For \texttt{env} well formed, $\Gamma'$ pointwise a type, and
\texttt{W : Ctx.Lift' l $\Gamma$ $\Gamma'$}: if the lifted term \texttt{e.lift' l} has a projection
translation over $\Gamma'$, then \texttt{e} has one (to some target) over $\Gamma$.
\end{theorem}
\begin{proof}
\uses{thm:isdefequ-weakn-iff}
\textbf{Not proved in Lean.} The declaration is \texttt{sorry}, as it was on master, but it now
carries a docstring saying why: the type of \texttt{e.lift' l} is only \emph{definitionally equal}
to a lifted type, so recovering the parameters over $\Gamma$ needs
\texttt{VEnv.IsDefEqU.weakN\_iff}, itself an admitted master lemma
(\srcloc{Lean4Lean/Theory/Typing/UniqueTyping.lean}{172}). The consequence is that
\texttt{TrExprS.weakFV'\_inv} (Theorem \ref{thm:trexprs-weakfv-inv}) and the three
\texttt{Conditionally*.weakN\_inv} lemmas remain sorry-dependent; they were on master too, so this
is not a regression.
\end{proof}
\axfoot{sorryAx: the declaration is itself sorry; blocked on VEnv.IsDefEqU.weakN\_iff (master)}

\begin{theorem}[The projection translation transports along defeq contexts and defeq sources]
\label{thm:trproj-defeqdfc}
\lean{Lean4Lean.TrProj.defeqDFC}
\leanok
\uses{def:trproj, fam:defeqdfc}
\contrib{trproj} \srcloc{Lean4Lean/Verify/Typing/Lemmas.lean}{749}
\thesisref{no direct counterpart; the model-side analogue is \texttt{IsDefEq.defeqDFC}}
For \texttt{env} well formed, a definitional equality of contexts
\texttt{IsDefEqCtx U [] $\Gamma_1$ $\Gamma_2$} and of sources
$\texttt{e}_1 \equiv \texttt{e}_2$ over $\Gamma_1$: a projection translation of \texttt{e}$_1$ over
$\Gamma_1$ yields one of \texttt{e}$_2$ over $\Gamma_2$, to a possibly different target.
\end{theorem}
\begin{proof}
\uses{fam:defeq-mixing}
\leanok
Four lines. The source occurs exactly once, as the argument of the expansion, whose own typing does
not mention it: transport \texttt{major\_ty} and \texttt{fn\_ty} along \texttt{HasType.defeqDFC},
replace the major by the definitionally equal one with \texttt{HasType.defeqU\_l}, and rebuild the
application. The statement deliberately produces only \emph{some} target; callers recover the
relation to a given target with the still-open \texttt{TrProj.uniq}.
\end{proof}
\axfoot{sorryAx via VEnv.HasType.defeqU\_l, i.e. unique typing (VEnv.IsDefEq.uniq), which rests on master's admitted injectivity lemmas (Theory/Typing/Injectivity.lean, e.g. IsDefEqU.forallE\_inv\_stratified) and, on this branch, additionally on VEnv.WF.patsStrong (Theory/Typing/EnvLemmas.lean:334, iota)}

\begin{theorem}[The projection translation is monotone in the environment]
\label{thm:trproj-mono}
\lean{Lean4Lean.TrProj.mono}
\leanok
\uses{def:trproj, struct:venv-le}
\contrib{trproj} \srcloc{Lean4Lean/Verify/Typing/Lemmas.lean}{767}
\thesisref{no counterpart}
If \texttt{env} $\le$ \texttt{env'} then \texttt{TrProj env U $\Gamma$ s i e e'} implies
\texttt{TrProj env' U $\Gamma$ s i e e'}.
\end{theorem}
\begin{proof}
\uses{thm:isdefeq-mono}
\leanok
Every field is monotone: \texttt{pat} by the \texttt{pats} component of \texttt{VEnv.LE},
\texttt{ctor} and \texttt{minor\_arity} by its \texttt{constants} component, and the two typings by
\texttt{HasType.mono}. The lemma has no master counterpart; it exists only because \texttt{TrProj}
now mentions \texttt{env}, and it rests on the \texttt{iota} branch's monotonicity of the pattern
registry, a neat cross-branch fit.
\end{proof}

\begin{theorem}[The expansion of a projection is well typed]
\label{thm:trproj-wf}
\lean{Lean4Lean.TrProj.wf}
\leanok
\uses{def:trproj}
\contrib{trproj} \srcloc{Lean4Lean/Verify/Typing/Lemmas.lean}{939}
\thesisref{Regularity (typesys.tex, \texttt{thm:reg}), at the level of the translation}
\texttt{TrProj env U $\Gamma$ s i e e'} implies \texttt{VExpr.WF env U $\Gamma$ e'}, with no
hypothesis on \texttt{e}.
\end{theorem}
\begin{proof}
\uses{def:hastype-istype}
\leanok
Unfold \texttt{e'} to the expansion applied to \texttt{e} and apply \texttt{fn\_ty} to
\texttt{major\_ty}. This is stronger than master's stub, which additionally assumed
\texttt{VExpr.WF env U $\Gamma$ e}: well-typedness of the major is now a field of the relation.
Master's stub was in fact mis-stated, its \texttt{TrProj} hypothesis being taken over an
auto-bound context variable unrelated to the $\Gamma$ of its conclusion. The
corresponding case of \texttt{TrExprS.wf} shrinks accordingly, from \texttt{h2.wf (ih h$\Delta$)} to
\texttt{h2.wf}.
\end{proof}

\begin{theorem}[Uniqueness of the projection translation (open)]
\label{thm:trproj-uniq}
\lean{Lean4Lean.TrProj.uniq}
\leanok
\uses{def:trproj, def:isdefeq}
\contrib{trproj} \srcloc{Lean4Lean/Verify/Typing/Lemmas.lean}{992}
\thesisref{Unique typing (unique.tex, \texttt{thm:utype}), at the level of the translation}
For \texttt{env} well formed and \texttt{IsDefEqCtx U [] $\Gamma_1$ $\Gamma_2$}: if
\texttt{TrProj env U $\Gamma_1$ s$_1$ i e$_1$ e$_1'$} and
\texttt{TrProj env U $\Gamma_2$ s$_2$ i e$_2$ e$_2'$} with
$\texttt{e}_1 \equiv \texttt{e}_2$ over $\Gamma_1$, then
$\texttt{e}_1' \equiv \texttt{e}_2'$ over $\Gamma_1$. Note that the two structure names are
separate variables: only the field index is shared, because the \texttt{proj} case of the model's
congruence compares two projections up to the index alone.
\end{theorem}
\begin{proof}
\uses{thm:unique-typing, def:proj-fn}
\textbf{Not proved in Lean.} The declaration is \texttt{sorry}, as on master, but it now carries an
eight-line docstring listing what a proof would need: unique typing of the projection function,
type-former injectivity for \texttt{usS}, \texttt{params} and the head names, and functionality of
the $\iota$ registry for \texttt{np} and the field count, since the key \texttt{np+1+1+0} does not
by itself exclude reading a two-minor recursor as \texttt{(np+1)+1+1+0} nor a second constructor for
\texttt{s}; then a congruence for \texttt{instPis} under pointwise definitional equality of
\texttt{params}, one for \texttt{projFn} by induction on \texttt{i}, and \texttt{appDF}.
\end{proof}
\axfoot{sorryAx: the declaration is itself sorry; TrExprS.uniq and TrExpr.proj inherit it}

\begin{remark}
The sketch does not say how the per-field level list \texttt{uss} is pinned. Two
\texttt{TrProjCtor} witnesses for the same projection may in principle choose different
\texttt{uss}, in which case the two expansions are syntactically different recursor applications at
different universe instances, and a definitional equality between them is not obviously available.
Either the sketch is incomplete or \texttt{TrProjCtor} needs a further field tying \texttt{uss i} to
the sort of the field type \texttt{F}$_i$.
\end{remark}

\begin{theorem}[Instantiation of the projection translation]
\label{thm:trproj-instn}
\lean{Lean4Lean.TrProj.instN}
\leanok
\uses{def:trproj, fam:ctx-lift}
\contrib{trproj} \srcloc{Lean4Lean/Verify/Typing/Lemmas.lean}{1296}
\thesisref{Substitution (typesys.tex, \texttt{thm:subst}), at the level of the translation}
If \texttt{env} is \texttt{Ordered}, \texttt{W : Ctx.InstN $\Gamma_0$ e$_0$ A$_0$ k $\Gamma_1$
$\Gamma$}, \texttt{TrProj env U $\Gamma_1$ s i e e'} and
$\Gamma_0 \vdash \texttt{e}_0 : \texttt{A}_0$, then
\texttt{TrProj env U $\Gamma$ s i (e.inst e$_0$ k) (e'.inst e$_0$ k)}.
\end{theorem}
\begin{proof}
\uses{fam:ctorheaded-stability, thm:proj-instpis-ctorheaded, fam:proj-fn-subst, fam:proj-motivebody-subst, fam:proj-instpis-subst, thm:isdefeq-instn, thm:isdefeq-closedn}
\leanok
The same record rebuild as \texttt{weak'}, with \texttt{params} and the field telescope instantiated
pointwise. The key ingredient is \texttt{piBinders\_inst\_of\_ctorHeaded}, which needs the
constructor type to be \texttt{CtorHeaded} (obtained from \texttt{instPis\_ctorHeaded}), since
instantiating inside a telescope could otherwise create new binders; the two typings go by
\texttt{HasType.instN} and the instantiation equation by \texttt{instPis\_inst} plus closedness of
the constant's recorded type. The \texttt{CtorHeaded} field of \texttt{TrProjCtor} exists precisely
for this lemma. Both the \texttt{Ordered} hypothesis and
$\Gamma_0 \vdash \texttt{e}_0 : \texttt{A}_0$ are new relative to master's stub; both are available
at the single call site inside \texttt{TrExprS.instN}
(\srcloc{Lean4Lean/Verify/Typing/Lemmas.lean}{1343}).
\end{proof}

\begin{theorem}[Level instantiation of the projection translation]
\label{thm:trproj-instl}
\lean{Lean4Lean.TrProj.instL}
\leanok
\uses{def:trproj, def:vlevel-inst}
\contrib{trproj} \srcloc{Lean4Lean/Verify/Typing/Lemmas.lean}{1588}
\thesisref{no separate counterpart; universe substitution is implicit in the thesis}
If every level of \texttt{ls} is well formed over \texttt{U'}, then
\texttt{TrProj env U $\Gamma$ s i e e'} implies
\texttt{TrProj env U' ($\Gamma$.map (VExpr.instL ls)) s i (e.instL ls) (e'.instL ls)}, with the
structure's level list, every per-field level list, the parameters and the field telescope
instantiated pointwise.
\end{theorem}
\begin{proof}
\uses{fam:pibinders-stability, fam:proj-fn-subst, fam:proj-motivebody-subst, fam:proj-instpis-subst, thm:isdefeq-instl}
\leanok
The same shape as \texttt{weak'} and \texttt{instN}, using \texttt{piBinders\_instL},
\texttt{projFn\_instL}, \texttt{projMotiveBody\_instL}, \texttt{HasType.instL} and
\texttt{instPis\_instL} with \texttt{instL\_instL}. This is the one place where the function-valued
\texttt{uss} really has to be transported, as
\texttt{fun j => (uss j).map (VLevel.inst ls)}, and it goes through cleanly, which is mild evidence
that the design is workable. Master's stub already carried \texttt{hls}, so unlike
\texttt{weak'} and \texttt{instN} this signature is unchanged.
\end{proof}

\section{Transport of the translation}

The lemmas of this section are the translation-level shadows of the thesis's structural lemmas.
All of them are master's. The \texttt{trproj} branch touched the \texttt{proj} case of five of the
inductions, \texttt{weakFV'}, \texttt{weakBV}, \texttt{mono}, \texttt{wf} and \texttt{instN}, plus
the signature of the \texttt{TrExpr.proj} builder, which is where the corresponding \texttt{TrProj}
lemma of the previous section is applied; the \texttt{proj} cases of \texttt{uniq},
\texttt{defeqDFC}, \texttt{weakFV'\_inv}, \texttt{instN\_let} and \texttt{abstract} needed no edit,
since the shape of the argument they pass did not change. The \texttt{iota} branch touched none of
them.

\begin{theorem}[Translation is preserved by free-variable weakening]
\label{thm:trexprs-weakfv}
\lean{Lean4Lean.TrExprS.weakFV', Lean4Lean.TrExpr.weakFV', Lean4Lean.TrExprS.weakFV, Lean4Lean.TrExpr.weakFV}
\leanok
\uses{ind:trexprs, ind:vt-fvlift, def:trexpr}
\contrib{mixed} \srcloc{Lean4Lean/Verify/Typing/Lemmas.lean}{669}
\thesisref{Weakening (typesys.tex, \texttt{thm:weak}), at the level of the translation}
For \texttt{env} \texttt{Ordered}, a lift \texttt{W : VLCtx.FVLift'} from $\Delta$ to $\Delta'$ and
$\Delta'$ well formed: \texttt{TrExprS env Us $\Delta$ e e'} implies
\texttt{TrExprS env Us $\Delta'$ e (e'.lift' (n.consN k))}, the source term being unchanged since
free variables are named. \texttt{weakFV} is the \texttt{FVLift}/\texttt{liftN} specialisation, and
both have \texttt{TrExpr} variants, which need \texttt{env} well formed because they also weaken
the definitional equality.
\end{theorem}
\begin{proof}
\uses{thm:trproj-weak, thm:isdefeq-weakn, family:vt-vlctx-wf-lemmas}
\leanok
Induction on the translation. Both variable cases go through \texttt{FVLift'.find?}; the
\texttt{app}, \texttt{lam}, \texttt{forallE} and \texttt{letE} rules weaken their typing side
conditions with \texttt{HasType.weak'} and recurse, extending the lift by \texttt{cons\_bvar} and
the context well-formedness by the weakened binder type at each binder. The \texttt{proj} case is
the only one the branch changed: it now passes the environment hypothesis on to
\texttt{TrProj.weak'} (Theorem \ref{thm:trproj-weak}).
\end{proof}

\begin{theorem}[Translation is preserved by bound-variable weakening]
\label{thm:trexprs-weakbv}
\lean{Lean4Lean.TrExprS.weakBV, Lean4Lean.TrExpr.weakBV}
\leanok
\uses{ind:trexprs, ind:vt-bvlift, def:trexpr}
\contrib{mixed} \srcloc{Lean4Lean/Verify/Typing/Lemmas.lean}{710}
\thesisref{Weakening (typesys.tex, \texttt{thm:weak}), at the level of the translation}
For \texttt{env} \texttt{Ordered} and a \texttt{VLCtx.BVLift} from $\Delta$ to $\Delta'$:
\texttt{TrExprS env Us $\Delta$ e e'} implies
\texttt{TrExprS env Us $\Delta'$ (e.liftLooseBVars' dk dn) (e'.liftN n k)}. Here the source term
does move, because the inserted entries are anonymous.
\end{theorem}
\begin{proof}
\uses{thm:trproj-weak, thm:isdefeq-weakn, family:vt-closed-lemmas}
\leanok
Induction on the translation, weakening each typing side condition with \texttt{HasType.weakN} and
each binder by \texttt{BVLift.cons}. The literal case is the only non-routine one: the constructor
form of a literal is closed, so \texttt{liftLooseBVars} is the identity on it, which is
\texttt{Closed.toConstructor.looseBVarRange\_le}. The \texttt{proj} case was updated on
\texttt{trproj} to pass the environment hypothesis to \texttt{TrProj.weakN}.
\end{proof}

\begin{theorem}[Translation is monotone in the environment]
\label{thm:trexprs-mono}
\lean{Lean4Lean.TrExprS.mono, Lean4Lean.TrExpr.mono, Lean4Lean.VEnv.ContainsLits.mono}
\leanok
\uses{ind:trexprs, struct:venv-le, def:vt-containslits}
\contrib{mixed} \srcloc{Lean4Lean/Verify/Typing/Lemmas.lean}{781}
\thesisref{no counterpart}
If \texttt{env} $\le$ \texttt{env'} then \texttt{TrExprS env Us $\Delta$ e e'} implies
\texttt{TrExprS env' Us $\Delta$ e e'}, and likewise for \texttt{TrExpr} and for the literal side
condition \texttt{ContainsLits}.
\end{theorem}
\begin{proof}
\uses{thm:trproj-mono, thm:isdefeq-mono}
\leanok
Induction on the translation: the \texttt{const} case uses the \texttt{constants} component of
\texttt{VEnv.LE} and the typing side conditions use \texttt{HasType.mono}. The \texttt{proj} case
changed on \texttt{trproj} from \texttt{exact .proj ih h2}, where there was nothing to do because
\texttt{TrProj} did not mention \texttt{env}, to \texttt{exact .proj ih (h2.mono henv)}; this is
the obligation that forced Theorem \ref{thm:trproj-mono} into existence.
\end{proof}

\begin{theorem}[A translation produces a well-formed model term]
\label{thm:trexprs-wf}
\lean{Lean4Lean.TrExprS.wf, Lean4Lean.TrExpr.wf, Lean4Lean.TrExprS.trExpr}
\leanok
\uses{ind:trexprs, def:vt-vlctx-wf, def:hastype-istype}
\contrib{mixed} \srcloc{Lean4Lean/Verify/Typing/Lemmas.lean}{947}
\thesisref{Regularity (typesys.tex, \texttt{thm:reg}), at the level of the translation}
Over an \texttt{Ordered} environment and a well-formed $\Delta$, \texttt{TrExprS env Us $\Delta$ e e'}
implies that \texttt{e'} has some type over $\Delta$\texttt{.toCtx}; \texttt{trExpr} packages that
into a \texttt{TrExpr} by reflexivity. The \texttt{TrExpr} form needs no hypotheses at all, since
the witness carries its own typing.
\end{theorem}
\begin{proof}
\uses{thm:trproj-wf, family:vt-vlctx-wf-lemmas, def:vlevel-oflevel}
\leanok
Induction on the translation, reading the typing off the side conditions of each rule: variables
use \texttt{VLCtx.WF.find?\_wf}, and the \texttt{sort} and \texttt{const} cases rebuild the model's
own rules from the level translation. On \texttt{trproj} the projection case simplified from an
application of the stub to the recursive call to a bare \texttt{h2.wf}, because \texttt{TrProj}
now carries the major premise's typing itself (Theorem \ref{thm:trproj-wf}).
\end{proof}

\begin{theorem}[Uniqueness of translation up to definitional equality]
\label{thm:trexprs-uniq}
\lean{Lean4Lean.TrExprS.uniq, Lean4Lean.TrExpr.uniq}
\leanok
\uses{ind:trexprs, ind:vt-vlctx-isdefeq, def:isdefeq}
For \texttt{env} well formed and $\Delta_1$, $\Delta_2$ definitionally equal translation contexts:
two translations of the \emph{same} source expression, one over $\Delta_1$ and one over $\Delta_2$,
have definitionally equal targets over $\Delta_1$\texttt{.toCtx}. This is the translation-level
form of Unique typing (unique.tex, \texttt{thm:utype}).
\end{theorem}
\begin{proof}
\uses{thm:trproj-uniq, fam:defeq-mixing, thm:trexprs-wf}
\leanok
Induction on the first translation, inverting the second at each node. The variable cases are
\texttt{VLCtx.IsDefEq.find?\_uniq}; \texttt{appDF}, \texttt{lamDF} and \texttt{forallEDF} assemble
the congruences, each preceded by an \texttt{of\_l} that re-types the recursive equality at the
left-hand side's type. The projection case delegates to \texttt{TrProj.uniq}
(Theorem \ref{thm:trproj-uniq}), which is \texttt{sorry}, so this theorem is sorry-dependent; it
already was on master.
\end{proof}
\axfoot{sorryAx via TrProj.uniq (Verify/Typing/Lemmas.lean:992)}

\begin{theorem}[Translation transports along definitionally equal contexts]
\label{thm:trexprs-defeqdfc}
\lean{Lean4Lean.TrExprS.defeqDFC, Lean4Lean.TrExprS.defeqDFC'}
\leanok
\uses{ind:trexprs, ind:vt-vlctx-isdefeq}
If \texttt{e} translates over $\Delta_1$ and $\Delta_1$ is definitionally equal to $\Delta_2$, then
\texttt{e} translates over $\Delta_2$ to \emph{some} target; the primed form packages the original
target back in as a \texttt{TrExpr} over $\Delta_2$.
\end{theorem}
\begin{proof}
\uses{thm:trexprs-uniq, thm:trproj-defeqdfc, fam:defeqdfc}
\leanok
Induction on the translation: at each node the typing side conditions are moved across with
\texttt{HasType.defeqDFC} and then repaired by \texttt{defeqU\_l} against the uniqueness of the
recursive results (Theorem \ref{thm:trexprs-uniq}); the projection case is
Theorem \ref{thm:trproj-defeqdfc}, which is why it produces only an existential target.
\end{proof}
\axfoot{sorryAx inherited from TrExprS.uniq, hence TrProj.uniq}

\begin{lemma}[Introduction rules for the up-to-defeq translation]
\label{family:trexpr-builders}
\lean{Lean4Lean.TrExpr.defeq, Lean4Lean.TrExpr.app, Lean4Lean.TrExpr.lam, Lean4Lean.TrExpr.forallE, Lean4Lean.TrExpr.letE, Lean4Lean.TrExpr.lit, Lean4Lean.TrExpr.mdata, Lean4Lean.TrExpr.proj}
\leanok
\uses{def:trexpr}
\contrib{mixed} \srcloc{Lean4Lean/Verify/Typing/Lemmas.lean}{967}
\thesisref{no counterpart}
One introduction rule per expression former, lifting the strict rules of \texttt{TrExprS} to
\texttt{TrExpr}, together with \texttt{TrExpr.defeq}, which post-composes a target with a
definitional equality. The binder rules carry the content: the body has been translated in a
context whose binder type is only definitionally equal to the one the rule is asked for.
\end{lemma}
\begin{proof}
\uses{thm:trexprs-uniq, thm:trexprs-defeqdfc}
\leanok
The three binder rules unpack their \texttt{TrExpr} witnesses, transport the body's strict
translation along the induced definitional equality of contexts with \texttt{TrExprS.defeqDFC},
realign the two targets with \texttt{TrExprS.uniq}, and rebuild the strict rule; \texttt{app},
\texttt{lit} and \texttt{mdata} only re-type the recursive equalities (\texttt{of\_r},
\texttt{appDF}) and rebuild. \texttt{TrExpr.proj}
(\srcloc{Lean4Lean/Verify/Typing/Lemmas.lean}{1162}) changed on \texttt{trproj} only in that its
\texttt{TrProj} argument gained the \texttt{env} and \texttt{Us.length} parameters; its proof
still routes through the admitted \texttt{TrProj.uniq}.
\end{proof}
\axfoot{sorryAx in defeq, app, lam, forallE, letE and proj, not in lit or mdata: defeq and app through the model's admitted unique typing alone, the binder rules and proj additionally through TrExprS.uniq, hence TrProj.uniq (Verify/Typing/Lemmas.lean:992)}

\begin{theorem}[Strengthening: a translation that does not use the extra variables descends]
\label{thm:trexprs-weakfv-inv}
\lean{Lean4Lean.TrExprS.weakFV'_inv, Lean4Lean.TrExprS.weakFV_inv}
\leanok
\uses{ind:trexprs, ind:vt-fvlift, def:vt-closed, def:vt-fvarsin}
For \texttt{env} well formed, a lift \texttt{W} from $\Delta$ to $\Delta_2$ and $\Delta_1$
definitionally equal to $\Delta_2$: if \texttt{e} translates over $\Delta_1$, is \texttt{Closed} at
the bound-variable depth and has all its free variables in $\Delta$\texttt{.fvars}, then \texttt{e}
translates over the smaller $\Delta$. This is the strengthening direction of Weakening
(typesys.tex, \texttt{thm:weak}, item 3), and the only place in the chapter where that direction
is needed.
\end{theorem}
\begin{proof}
\uses{thm:trproj-weak-inv, thm:trexprs-uniq, thm:trexprs-weakfv, fam:weakening-iff}
\leanok
Induction on the translation. At each node the recursive results are weakened back up with
Theorem \ref{thm:trexprs-weakfv}, realigned with Theorem \ref{thm:trexprs-uniq}, and the typing
side conditions then pushed down by \texttt{VExpr.WF.weak'\_iff} and \texttt{HasType.weak'\_iff},
both consequences of unique typing. The bound-variable case is an inner induction on the lift.
The projection case calls the admitted \texttt{TrProj.weak'\_inv}
(Theorem \ref{thm:trproj-weak-inv}); that is the reason this theorem, and with it the three
\texttt{Conditionally} strengthening lemmas of the last section of this chapter, stay
sorry-dependent.
\end{proof}
\axfoot{sorryAx via TrProj.weak'\_inv, itself blocked on VEnv.IsDefEqU.weakN\_iff (Theory/Typing/UniqueTyping.lean:172, master)}

\begin{theorem}[Translation commutes with instantiating a lambda binder]
\label{thm:trexprs-instn}
\lean{Lean4Lean.TrExprS.instN_var, Lean4Lean.TrExprS.instN, Lean4Lean.TrExprS.inst, Lean4Lean.TrExpr.inst}
\leanok
\uses{ind:trexprs, ind:vt-instn}
\contrib{mixed} \srcloc{Lean4Lean/Verify/Typing/Lemmas.lean}{1251}
\thesisref{Substitution (typesys.tex, \texttt{thm:subst}), at the level of the translation}
Let \texttt{env} be \texttt{Ordered}, let \texttt{e0} translate to \texttt{e0'} over $\Delta_0$
with $\Delta_0 \vdash \texttt{e0'} : \texttt{A0}$, and let \texttt{W} be a \texttt{VLCtx.InstN}
substituting \texttt{e0'} for the binder at depth \texttt{dk} of $\Delta_1$, giving $\Delta$. Then
\texttt{TrExprS env Us $\Delta_1$ e e'} implies
\texttt{TrExprS env Us $\Delta$ (e.instantiate1' e0 dk) (e'.inst e0' k)}. \texttt{inst} is the
outermost-binder case and \texttt{TrExpr.inst} its up-to-defeq form.
\end{theorem}
\begin{proof}
\uses{thm:trproj-instn, thm:isdefeq-instn, thm:trexprs-weakbv}
\leanok
The variable case is a separate induction on \texttt{W} (\texttt{instN\_var}), splitting on whether
the variable is the one being substituted, a bound variable inside or outside the substituted
position, or a free variable; when the substituted value has to cross a binder it is re-lifted by
Theorem \ref{thm:trexprs-weakbv}. The rest is an induction on the translation, instantiating each
typing side condition by \texttt{HasType.instN}, with the literal case again using closedness of
the constructor form. The projection case reads \texttt{h2.instN henv W.toCtx t0} on
\texttt{trproj}, where master's stub took no typing argument; the new hypothesis is available at
this single call site.
\end{proof}
\axfoot{sorryAx: the TrExpr form routes through the TrExpr builders and TrExprS.uniq, hence TrProj.uniq}

\begin{theorem}[Translation commutes with unfolding a let binder]
\label{thm:trexprs-instlet}
\lean{Lean4Lean.TrExprS.instN_let_var, Lean4Lean.TrExprS.instN_let, Lean4Lean.TrExprS.inst_let, Lean4Lean.TrExpr.inst_let}
\leanok
\uses{ind:trexprs, ind:vt-instn}
The analogue of Theorem \ref{thm:trexprs-instn} for a \texttt{vlet} entry: the source term is
substituted but the model term and the model context are unchanged, since a \texttt{vlet} occupies
no slot in \texttt{toCtx}. The thesis's counterpart is $\zeta$ reduction (axioms.tex, let binders).
\end{theorem}
\begin{proof}
\uses{thm:trexprs-weakbv}
\leanok
The same two-step shape as Theorem \ref{thm:trexprs-instn}: a variable lemma by induction on
\texttt{InstLet}, then an induction on the translation. Because \texttt{InstLet.toCtx} says the
model context does not move, no typing side condition has to be transported at all.
\end{proof}
\axfoot{sorryAx: the TrExpr form routes through the TrExpr builders and TrExprS.uniq}

\begin{theorem}[Translation commutes with abstracting a free variable]
\label{thm:trexprs-abstract}
\lean{Lean4Lean.TrExprS.abstract, Lean4Lean.TrExpr.abstract}
\leanok
\uses{ind:trexprs, ind:vt-abstract}
For an \texttt{Abstract} relation turning the named entry \texttt{v0} of $\Delta_1$ into an
anonymous one, giving $\Delta$: if \texttt{e} translates over $\Delta_1$ to \texttt{e'}, then
\texttt{e.abstract1 v0 dk} translates over $\Delta$ to the \emph{same} \texttt{e'}. The model side
does not move, which is the point of the relation.
\end{theorem}
\begin{proof}
\uses{family:vt-fvarsin-lemmas}
\leanok
Induction on the translation; the two variable cases carry the content.
\texttt{Abstract.find?} handles every variable other than \texttt{v0}, and
\texttt{Abstract.find?\_self} turns an occurrence of \texttt{v0} itself into the new bound
variable. The literal case uses \texttt{FVarsIn.abstract\_eq\_self}, and the typing side conditions
are transported by a rewrite along \texttt{Abstract.toCtx}.
\end{proof}

\begin{lemma}[Opening and closing binders with fresh free variables]
\label{family:trexprs-uninstantiate}
\lean{Lean4Lean.TrExprS.uninstantiateN, Lean4Lean.TrExpr.uninstantiateN, Lean4Lean.TrExprS.uninstantiate, Lean4Lean.TrExpr.uninstantiate, Lean4Lean.TrExprS.inst_fvar, Lean4Lean.TrExpr.rebuild_mkAppRevList, Lean4Lean.TrExpr.rebuild_mkAppList}
\leanok
\uses{ind:trexprs, ind:vt-abstract, def:mkapps}
Translating a body that has been opened with a fresh free variable is interchangeable with
translating it under an anonymous binder: \texttt{uninstantiate} closes the binder again, provided
the body does not already mention that variable, and \texttt{inst\_fvar} is the opposite
direction. \texttt{rebuild\_mkAppList} re-applies an already translated argument spine to a
definitionally equal head. These are the lemmas the checker uses around each
\texttt{withLocalDecl}.
\end{lemma}
\begin{proof}
\uses{thm:trexprs-abstract, thm:trexprs-instn, thm:trexprs-weakfv}
\leanok
The \texttt{uninstantiate} family is Theorem \ref{thm:trexprs-abstract} composed with
\texttt{FVarsIn.abstract\_instantiate1}, which cancels an abstraction against an instantiation at a
variable not occurring in the term. \texttt{inst\_fvar} weakens past the new named entry and then
instantiates, splitting on \texttt{vlam} versus \texttt{vlet}.
\texttt{rebuild\_mkAppRevList} is an induction on the argument list using the \texttt{app} builder
of Lemma \ref{family:trexpr-builders}, with \texttt{TrExprS.uniq} closing the empty case, and
\texttt{rebuild\_mkAppList} is a one-line rewrite through \texttt{mkAppRevList\_reverse}.
\end{proof}
\axfoot{sorryAx in the two rebuild lemmas only, via the TrExpr app builder and TrExprS.uniq, hence TrProj.uniq; the uninstantiate family itself is clean}

\begin{definition}[Syntactic (on-the-nose) uniqueness of translation]
\label{def:trexprs-isunique}
\lean{Lean4Lean.TrExprS.IsUnique, Lean4Lean.TrExprS.IsUnique.natLitToConstructor, Lean4Lean.TrExprS.IsUnique.strLitToConstructor, Lean4Lean.TrExprS.IsUnique.toConstructor, Lean4Lean.TrExprS.IsUniqueDecl, Lean4Lean.TrExprS.IsUniqueCtx, Lean4Lean.TrExprS.IsUniqueCtx.find?_uniq, Lean4Lean.TrExprS.unique', Lean4Lean.TrExprS.unique}
\leanok
\uses{ind:trexprs}
\texttt{IsUnique e} is the structural predicate marking the expressions whose translation is
determined up to \emph{syntactic} equality. Every clause is \texttt{True} or the conjunction of
the subterms' clauses, with one exception: the projection clause is \texttt{False}.
\texttt{IsUniqueCtx} relates two contexts differing only in the \emph{types} of their declarations,
and \texttt{TrExprS.unique'} says that over such a pair an \texttt{IsUnique} expression has
literally equal targets. The proof is an induction on the first translation inverting the second,
the variable cases being \texttt{IsUniqueCtx.find?\_uniq}.
\end{definition}

\begin{remark}
The \texttt{proj} clause being \texttt{False} is where the absence of a projection node in
\texttt{VExpr} becomes visible: a projection has no canonical target, only a
definitional-equality class. The \texttt{trproj} work does not change this and arguably could
not, since \texttt{TrProjCtor} leaves the per-field level list \texttt{uss} free
(Definition \ref{struct:trprojctor}).
\end{remark}

\section{Universe substitution}

\begin{lemma}[Level translation and the smart level constructors]
\label{family:vt-oflevel}
\lean{Lean4Lean.ofLevel_mkLevelMax', Lean4Lean.ofLevel_isNeverZero, Lean4Lean.ofLevel_isAlwaysZero, Lean4Lean.ofLevel_mkLevelIMax', Lean4Lean.substParams_wf, Lean4Lean.substParams_wf_list}
\leanok
\uses{def:vlevel-oflevel, def:vlevel-le-equiv}
\texttt{VLevel.ofLevel} commutes, up to the model's level equivalence, with Lean's normalising
constructors \texttt{mkLevelMax'} and \texttt{mkLevelIMax'}, and it respects the
\texttt{isNeverZero} and \texttt{isAlwaysZero} tests those constructors branch on.
\texttt{substParams\_wf} lifts this to universe-parameter substitution: if \texttt{u} translates
over the parameter list \texttt{ps}, then \texttt{u.substParams'} translates over \texttt{Us} to
something equivalent to the instantiated translation. \texttt{substParams\_wf\_list} is the list
form. Equivalence rather than equality is unavoidable, since the real constructors normalise.
\end{lemma}
\begin{proof}
\uses{fam:vlevel-lattice}
\leanok
Case analysis on the level shapes, every branch of the real implementation's \texttt{subsumes}
test discharged by a lattice fact about \texttt{VLevel} (\texttt{le\_max\_left},
\texttt{max\_self}, \texttt{succ\_le\_succ}). \texttt{substParams\_wf} is then an induction on the
level whose parameter case reduces to a lookup in the substitution list, the \texttt{max} and
\texttt{imax} cases splitting on whether the smart constructor fired.
\end{proof}

\begin{theorem}[Translation commutes with universe-parameter substitution]
\label{thm:trexprs-instl}
\lean{Lean4Lean.TrExprS.instL, Lean4Lean.TrExpr.instL}
\leanok
\uses{ind:trexprs, def:instl, def:trexpr}
For \texttt{env} well formed and $\Delta$ well formed, if the levels \texttt{ls} translate to
\texttt{ls'} over \texttt{Us} and \texttt{ps.length = ls.length}, then a
translation over the parameter list \texttt{ps} and $\Delta$ yields
\texttt{TrExpr env Us ($\Delta$.instL ls') (e.instantiateLevelParams ps ls) (e'.instL ls')}. The
conclusion is a \texttt{TrExpr} and not a \texttt{TrExprS}: because Lean's level substitution
normalises, the substituted source translates only to something definitionally equal to the
substituted target.
\end{theorem}
\begin{proof}
\uses{family:vt-oflevel, thm:trproj-instl, thm:isdefeq-instl, family:trexpr-builders}
\leanok
Induction on the translation. The \texttt{sort} and \texttt{const} cases are where the equivalence
enters, through \texttt{substParams\_wf} and \texttt{substParams\_wf\_list} followed by the model's
\texttt{sortDF} and \texttt{constDF}; every other case applies the corresponding \texttt{TrExpr}
builder of Lemma \ref{family:trexpr-builders} to the level-instantiated side conditions. The
projection case is Theorem \ref{thm:trproj-instl}.
\end{proof}
\axfoot{sorryAx via the TrExpr builders and TrExprS.uniq, hence TrProj.uniq}

\section{Scopes, beta reduction and spines}

\begin{definition}[Upward-closed sets of free variables, and the typing bundle]
\label{def:isfvarupset}
\lean{Lean4Lean.VLocalDecl.ClosedN, Lean4Lean.VLCtx.Closed, Lean4Lean.IsFVarUpSet, Lean4Lean.IsFVarUpSet.congr, Lean4Lean.IsFVarUpSet.and_fvars, Lean4Lean.IsFVarUpSet.trivial, Lean4Lean.IsFVarUpSet.fvars, Lean4Lean.AllAbove, Lean4Lean.AllAbove.wf, Lean4Lean.FVarsBelow, Lean4Lean.FVarsBelow.rfl, Lean4Lean.FVarsBelow.trans, Lean4Lean.TrTyping}
\leanok
\uses{def:vt-fvarsin, def:vt-vlctx, ind:trexprs}
\texttt{IsFVarUpSet P $\Delta$} says that \texttt{P} is closed under the dependency edges recorded
in $\Delta$: if it holds of a named entry, it holds of every variable that entry's declaration
depends on. \texttt{FVarsBelow $\Delta$ e e'} says that \texttt{e'} satisfies every up-set that
\texttt{e} satisfies; this is the precise form of ``the computed type does not escape the scope of
the term''. \texttt{AllAbove} is the relativisation that makes an arbitrary predicate into one,
and \texttt{TrTyping} bundles \texttt{FVarsBelow} with the translations of a term and of its type
and the model typing relating them. It is the postcondition of the checker's \texttt{inferType},
and the syntactic half of Regularity (typesys.tex, \texttt{thm:reg}).
\texttt{VLCtx.Closed} and \texttt{VLocalDecl.ClosedN} are the auxiliary closedness predicates on
the context itself.
\end{definition}

\begin{definition}[Source-level beta reduction and \texttt{cheapBetaReduce}]
\label{ind:betareduce}
\lean{Lean4Lean.LambdaBodyN, Lean4Lean.LambdaBodyN.closed, Lean4Lean.LambdaBodyN.add, Lean4Lean.LambdaBodyN.instantiateList, Lean4Lean.LambdaBodyN.instantiateRevList, Lean4Lean.BetaReduce, Lean4Lean.BetaReduce.mkAppRevList, Lean4Lean.BetaReduce.mkAppList, Lean4Lean.BetaReduce.instantiateList, Lean4Lean.BetaReduce.instantiateRevList, Lean4Lean.FVarsBelow.betaReduce, Lean4Lean.BetaReduce.inst_reduce, Lean4Lean.BetaReduce.cheapBetaReduce, Lean4Lean.FVarsBelow.cheapBetaReduce}
\leanok
\uses{def:isfvarupset, def:kernel-cheap-beta, family:expr-subst}
\texttt{LambdaBodyN n e b} peels \texttt{n} lambda binders off \texttt{e}; \texttt{BetaReduce} is
the reflexive-transitive congruence closure, along the function side of an application, of a
single $\beta$ step whose argument is closed (\texttt{looseBVarRange' = 0}). Two results matter:
\texttt{BetaReduce.cheapBetaReduce} says that Lean's substitution-free
\texttt{Expr.cheapBetaReduce} is such a reduction on closed terms, and
\texttt{FVarsBelow.betaReduce} says that $\beta$ never enlarges the free-variable set. The first is
one of the larger master proofs in the file: the real implementation's \texttt{loop} is analysed
by an invariant producing a \texttt{LambdaBodyN} witness, after which its \texttt{cont} branches
are matched against an \texttt{inst\_reduce} of the spine. Together they let the checker's
\texttt{whnfCore} call \texttt{cheapBetaReduce} without losing its specification. The thesis's
counterpart is the $\beta$ rule (axioms.tex, Reduction).
\end{definition}

\begin{theorem}[Beta reduction preserves the translation]
\label{thm:trexpr-beta}
\lean{Lean4Lean.TrExpr.beta, Lean4Lean.TrExpr.cheapBetaReduce}
\leanok
\uses{ind:betareduce, def:trexpr}
For \texttt{env} well formed and $\Delta$ well formed: if \texttt{TrExpr env Us $\Delta$ e e'}
and \texttt{BetaReduce e e2}, then \texttt{TrExpr env Us $\Delta$ e2 e'}: the model target is unchanged, the source reduct being
definitionally equal on the nose. \texttt{cheapBetaReduce} is the instance the checker uses, with
the closedness hypothesis supplied by Theorem \ref{thm:trexprs-closed-fvars} together with
$\Delta$\texttt{.NoBV}.
\end{theorem}
\begin{proof}
\uses{thm:trexprs-instn, thm:trexprs-uniq, thm:trexprs-defeqdfc, thm:isdefeq-instdf, thm:isdefeq-foralle-inv, fam:strong-inversions, family:trexpr-builders}
\leanok
Induction on the beta derivation. The congruence case is the \texttt{app} builder; the base case
inverts the application and the lambda, recovers the binder type with \texttt{lam\_inv} and
\texttt{forallE\_inv}, realigns the body's translation across the resulting definitional equality
of contexts with Theorems \ref{thm:trexprs-defeqdfc} and \ref{thm:trexprs-uniq}, and closes with
the model's \texttt{beta} rule and \texttt{instDF}.
\end{proof}
\axfoot{sorryAx via TrExprS.uniq, hence TrProj.uniq}

\begin{lemma}[Transport along Lean's structural equality]
\label{family:trexprs-eqv}
\lean{Lean4Lean.TrExprS.eqv, Lean4Lean.TrExpr.eqv, Lean4Lean.fvarsList_eqv, Lean4Lean.FVarsIn.eqv, Lean4Lean.FVarsBelow.eqv, Lean4Lean.TrTyping.eqv}
\leanok
\uses{ind:trexprs, def:isfvarupset, def:vt-fvarsin}
Everything this layer states is invariant under \texttt{Expr.eqv'}, the relation underlying
\texttt{==}, which differs from propositional equality only in binder names and
\texttt{BinderInfo}: the two translation relations, \texttt{fvarsList}, \texttt{FVarsIn},
\texttt{FVarsBelow} and the \texttt{TrTyping} bundle all transport along it. This is what makes
the checker's use of \texttt{==} on expressions sound at the specification level.
\end{lemma}
\begin{proof}
\uses{family:expr-beq}
\leanok
Simultaneous induction on the two expressions, discarding the mismatched shapes by
\texttt{rintro}, unfolding \texttt{Expr.eqv'} and closing the remaining cases with
\texttt{grind}.
\end{proof}

\begin{definition}[Application spines with their typing data]
\label{ind:appstack}
\lean{Lean4Lean.AppStack, Lean4Lean.AppStack.build_rev, Lean4Lean.AppStack.tr, Lean4Lean.AppStack.append, Lean4Lean.AppStack.build}
\leanok
\uses{ind:trexprs, def:mkapps}
\texttt{AppStack env Us $\Delta$ f f' as} records that the head \texttt{f} translates to
\texttt{f'} and that applying it to the arguments \texttt{as} one at a time stays translatable,
keeping the function and argument typings of every partial application along the way.
\texttt{build} produces one from a translation of the saturated spine \texttt{f.mkAppList as}, by
an induction on the argument list through \texttt{build\_rev}; \texttt{tr} projects the head's
translation back out.
\end{definition}

\begin{theorem}[Inversion for the translation of an application spine]
\label{thm:trexprs-mkapplist-inv}
\lean{Lean4Lean.TrExprS.mkAppList_inv}
\leanok
\uses{ind:trexprs, def:mkapps}
\contrib{trproj} \srcloc{Lean4Lean/Verify/Typing/Lemmas.lean}{2349}
\thesisref{no counterpart; implementation-level inversion}
If \texttt{f.mkAppList as} translates to \texttt{e'}, then there are \texttt{f'} and \texttt{as'}
with \texttt{TrExprS env Us $\Delta$ f f'}, a pointwise \texttt{List.Forall}$_2$ of translations
between \texttt{as} and \texttt{as'}, and \texttt{e' = f'.mkApps as'}.
\end{theorem}
\begin{proof}
\leanok
Structural recursion on \texttt{as}. The recursive call at head \texttt{.app f a} returns the
decomposition of the shorter spine, and inverting the \texttt{app} rule of \texttt{TrExprS} peels
the outermost argument off. Eight lines, with no side conditions.
\end{proof}

\begin{remark}
It is used five times in the projection and $\iota$ WHNF proof,
\srcloc{Lean4Lean/Verify/TypeChecker/WHNF.lean}{87} to \texttt{:93}. It partly duplicates
\texttt{AppStack.build} (Definition \ref{ind:appstack}), which sits immediately above it at
\srcloc{Lean4Lean/Verify/Typing/Lemmas.lean}{2344} and gives the same decomposition in a different
shape; the pointwise-list form is the one the WHNF proof needs, so the duplication is defensible,
but it leaves two spine-inversion idioms side by side.
\end{remark}

\section{Literals and the primitive invariant}

\begin{definition}[Model-side spellings of the literal constants]
\label{family:vt-model-literals}
\lean{Lean4Lean.VExpr.bool, Lean4Lean.VExpr.boolTrue, Lean4Lean.VExpr.boolFalse, Lean4Lean.VExpr.boolLit, Lean4Lean.VExpr.nat, Lean4Lean.VExpr.natZero, Lean4Lean.VExpr.natSucc, Lean4Lean.VExpr.natLit, Lean4Lean.VExpr.char, Lean4Lean.VExpr.string, Lean4Lean.VExpr.stringOfList, Lean4Lean.VExpr.listChar, Lean4Lean.VExpr.listCharNil, Lean4Lean.VExpr.listCharCons, Lean4Lean.VExpr.charOfNat, Lean4Lean.VExpr.listCharLit, Lean4Lean.VExpr.trLiteral, Lean4Lean.VExpr.boolOp2, Lean4Lean.VExpr.natOp2, Lean4Lean.VExpr.inst_natLit, Lean4Lean.VExpr.liftN_natLit, Lean4Lean.VExpr.insts_natLit, Lean4Lean.VExpr.subst_natLit, Lean4Lean.VExpr.inst_boolLit, Lean4Lean.VExpr.insts_boolLit, Lean4Lean.VExpr.subst_boolLit}
\leanok
\uses{ind:vexpr, def:vt-containslits}
Abbreviations naming the model terms the literal translation targets: \texttt{bool} with its two
constructors and \texttt{boolLit}; \texttt{nat}, \texttt{natZero}, \texttt{natSucc} and the
\emph{unary} numeral \texttt{natLit n}; \texttt{char}, \texttt{string}, \texttt{listChar} with its
two constructors, \texttt{charOfNat}, \texttt{stringOfList} and the character-list literal; and
\texttt{trLiteral}, which sends a \texttt{Lean.Literal} to the corresponding model term.
\texttt{boolOp2} and \texttt{natOp2} name the two halves of \texttt{Nat.bitwise}'s recorded type.
Alongside them are the \texttt{simp} lemmas saying that lifting, single and telescope
instantiation, and general substitution leave these closed terms alone; each is \texttt{rfl} or a
one-line induction on the numeral.
\end{definition}

\begin{definition}[Reflection of a metatheoretic function by a model term]
\label{family:vt-reflects}
\lean{Lean4Lean.VEnv.ReflectsNatNat', Lean4Lean.VEnv.ReflectsNatNatNat', Lean4Lean.VEnv.ReflectsNatNatBool', Lean4Lean.VEnv.ReflectsBoolBoolBool', Lean4Lean.VEnv.ReflectsNatNat, Lean4Lean.VEnv.ReflectsNatNatNat, Lean4Lean.VEnv.ReflectsNatNatBool}
\leanok
\uses{family:vt-model-literals, def:hastype-istype}
A term \texttt{e} \emph{reflects} a function \texttt{f} when two things hold: \texttt{e} has the
corresponding arrow type in the empty context, and applied to literals it is definitionally equal
to the literal answer. The primed forms speak of an arbitrary term, so that a partial application
such as \texttt{Nat.bitwise f} is covered; the unprimed forms specialise to \texttt{const fc []}
and are guarded by the constant's presence in the environment.
\end{definition}

\begin{remark}
The docstrings explain carefully why the typing clause cannot be dropped: a \texttt{Nat} primitive
is consumed as an \emph{operator} inside another primitive's recurrence, and congruence under an
application needs the codomain to be known non-dependent, which inverting an application does not
give. For \texttt{ReflectsBoolBoolBool'} there is a second reason, spelled out in the source: a
term of type \texttt{(x : Bool)} to \texttt{(if x then Bool else Bool)} satisfies the evaluation
clause but is not a \texttt{boolOp2}. This is good documentation of a non-obvious design choice.
\end{remark}

\begin{definition}[Reflection specification for \texttt{Nat.bitwise}]
\label{def:vt-reflects-bitwise}
\lean{Lean4Lean.VEnv.ReflectsNatBitwise}
\leanok
\uses{family:vt-reflects}
\texttt{Nat.bitwise} must be declared with exactly the recorded type
\texttt{boolOp2} to \texttt{natOp2}, and for \emph{every} larger well-formed environment and every
term reflecting a boolean binary operation \texttt{g}, the application \texttt{Nat.bitwise f} must
reflect \texttt{Nat.bitwise g}. The quantification over the larger environment is load-bearing and
documented as such: the hypothesis on \texttt{f} is an \texttt{IsDefEqU}, which is monotone, so in
hypothesis position it makes the statement anti-monotone; without the quantifier the fact could
not be carried from the environment where \texttt{Nat.bitwise} is checked to the one where
\texttt{Nat.land} supplies its operator.
\end{definition}

\begin{definition}[The primitive invariant as a name-keyed table of specification shapes]
\label{ind:primspec}
\lean{Lean4Lean.PrimSpec, Lean4Lean.PrimSpec.Holds, Lean4Lean.primSpecs, Lean4Lean.VEnv.HasPrimitives}
\leanok
\uses{family:vt-reflects, def:vt-reflects-bitwise, def:kernel-primitives}
\texttt{PrimSpec} enumerates the seven shapes a primitive specification can take: presence implies
presence, exact recorded type, the three \texttt{Reflects} shapes, the bitwise shape, and the
\texttt{String.ofList} shape. \texttt{PrimSpec.Holds env n} interprets one at a name;
\texttt{primSpecs} is the table of twenty-four entries, keyed by name in the same order as
\texttt{Environment.primitives}; and \texttt{VEnv.HasPrimitives env} says every entry of the table
holds. The design rationale is in the section docstring: each spec looks up exactly one name and
everything else it asserts is monotone in the environment, so environment extension becomes a
single theorem whose case analysis ranges over seven shapes rather than twenty-odd primitives.
\end{definition}

\begin{lemma}[Named accessors for the primitive table]
\label{family:hasprimitives-accessors}
\lean{Lean4Lean.VEnv.HasPrimitives.bool, Lean4Lean.VEnv.HasPrimitives.nat, Lean4Lean.VEnv.HasPrimitives.boolFalse, Lean4Lean.VEnv.HasPrimitives.boolTrue, Lean4Lean.VEnv.HasPrimitives.natZero, Lean4Lean.VEnv.HasPrimitives.natSucc, Lean4Lean.VEnv.HasPrimitives.charOfNat, Lean4Lean.VEnv.HasPrimitives.stringOfList, Lean4Lean.VEnv.HasPrimitives.natPred, Lean4Lean.VEnv.HasPrimitives.natAdd, Lean4Lean.VEnv.HasPrimitives.natSub, Lean4Lean.VEnv.HasPrimitives.natMul, Lean4Lean.VEnv.HasPrimitives.natPow, Lean4Lean.VEnv.HasPrimitives.natGcd, Lean4Lean.VEnv.HasPrimitives.natMod, Lean4Lean.VEnv.HasPrimitives.natDiv, Lean4Lean.VEnv.HasPrimitives.natBEq, Lean4Lean.VEnv.HasPrimitives.natBLE, Lean4Lean.VEnv.HasPrimitives.natBitwise, Lean4Lean.VEnv.HasPrimitives.natLAnd, Lean4Lean.VEnv.HasPrimitives.natLOr, Lean4Lean.VEnv.HasPrimitives.natXor, Lean4Lean.VEnv.HasPrimitives.natShiftLeft, Lean4Lean.VEnv.HasPrimitives.natShiftRight}
\leanok
\uses{ind:primspec}
One projection per entry of \texttt{primSpecs}, so that call sites keep naming the primitive rather
than the table index.
\end{lemma}
\begin{proof}
\leanok
Each is the invariant applied to the entry, with membership discharged by
\texttt{simp [primSpecs]}; the two \texttt{containsImplies} entries additionally project the two
names out of the resulting list.
\end{proof}

\begin{lemma}[Translation and typing of the literal constants]
\label{family:trexprs-literals}
\lean{Lean4Lean.TrExprS.boolFalse, Lean4Lean.TrExprS.boolTrue, Lean4Lean.TrExprS.boolLit, Lean4Lean.VExpr.WF.boolLit_has_type, Lean4Lean.TrExprS.lit_has_type, Lean4Lean.TrExprS.nat_of_natZero, Lean4Lean.TrExprS.natZero, Lean4Lean.TrExprS.natSucc, Lean4Lean.TrExprS.natLit, Lean4Lean.TrExprS.stringOfList, Lean4Lean.TrExprS.charOfNat, Lean4Lean.VEnv.HasPrimitives.nat_of_charOfNat, Lean4Lean.TrExprS.listChar, Lean4Lean.TrExprS.listCharNil, Lean4Lean.TrExprS.listCharCons, Lean4Lean.TrExprS.listCharLit, Lean4Lean.TrExprS.trLiteral, Lean4Lean.VExpr.instL_boolFalse, Lean4Lean.VExpr.instL_boolTrue, Lean4Lean.VExpr.instL_boolLit, Lean4Lean.VExpr.instL_natZero, Lean4Lean.VExpr.instL_natSucc, Lean4Lean.VExpr.instL_natLit, Lean4Lean.FVarsIn.boolLit}
\leanok
\uses{ind:trexprs, ind:primspec, family:vt-model-literals}
\contrib{mixed} \srcloc{Lean4Lean/Verify/Typing/Lemmas.lean}{1824}
\thesisref{no counterpart; the thesis does not model literals}
Given the primitive invariant and the presence of the relevant constants, each literal-related
source expression translates to the model term of
Definition \ref{family:vt-model-literals} and is well typed at the expected type: the two
\texttt{Bool} constructors and \texttt{boolLit}, the \texttt{Nat} numerals,
\texttt{Char.ofNat} and the character lists, \texttt{String.ofList}, and
\texttt{Literal.toConstructor} in general. Two lemmas run the other way,
recovering the presence of \texttt{Nat} or \texttt{Bool} from a well-typed literal.
\end{lemma}
\begin{proof}
\uses{fam:strong-inversions, struct:orderedstrong, thm:isdefeq-instl, def:kernel-lit-to-ctor}
\leanok
Each unfolds the relevant \texttt{HasPrimitives} accessor to pin the constant's recorded type,
applies the \texttt{const} rule of \texttt{TrExprS}, and moves the closed typing into the ambient
level and local context with \texttt{instL} and \texttt{weak0}; the character-list and
\texttt{trLiteral} lemmas are then inductions over the string. The reverse lemmas invert the
typing with \texttt{HasType.const\_inv}.
\end{proof}

\begin{remark}
Eight signature lines here were changed on the \texttt{iota} branch from \texttt{env.Ordered} to
\texttt{env.OrderedStrong} (\srcloc{Lean4Lean/Verify/Typing/Lemmas.lean}{1856} onwards), together
with two call sites where \texttt{wf.constWF} became \texttt{wf.ordered.constWF}. The change is
forced by \texttt{HasType.const\_inv} alone
(\srcloc{Lean4Lean/Theory/Typing/Strong.lean}{915}), whose own hypothesis moved because the
strengthening lemma it invokes, \texttt{IsDefEq.strong}
(\srcloc{Lean4Lean/Theory/Typing/Strong.lean}{805}), now needs subject reduction of the registered
$\iota$ rules. The other inversion used here, \texttt{HasType.forallE\_inv}
(\srcloc{Lean4Lean/Theory/Typing/Lemmas.lean}{854}), still takes only \texttt{Ordered} and is
simply fed the stronger hypothesis. The edit is mechanical, but it means every \emph{use} of these
lemmas
(\srcloc{Lean4Lean/Verify/TypeChecker/Reduce.lean}{16},
\srcloc{Lean4Lean/Verify/TypeChecker/IsDefEq.lean}{406},
\srcloc{Lean4Lean/Verify/TypeChecker/InferType.lean}{420}) now discharges \texttt{OrderedStrong}
through \texttt{VEnv.WF.orderedStrong}, which rests on the admitted \texttt{VEnv.WF.patsStrong}
(\srcloc{Lean4Lean/Theory/Typing/EnvLemmas.lean}{334}). A hypothesis strengthening in this group
therefore imports a new admission into the checker's literal handling.
\end{remark}

\begin{theorem}[A primitive specification survives environment extension]
\label{thm:primspec-addconst}
\lean{Lean4Lean.VEnv.addConst_eq_of_ne, Lean4Lean.PrimSpec.Holds.addConst, Lean4Lean.PrimSpec.Holds.addDefEq}
\leanok
\uses{ind:primspec, def:venv-basic, struct:venv-le}
Adding a constant under a name different from the spec's key, or adding a definitional equation,
preserves \texttt{PrimSpec.Holds}. \texttt{addConst\_eq\_of\_ne} is the supporting fact that a
successful \texttt{addConst} does not disturb the lookup of any other name.
\end{theorem}
\begin{proof}
\uses{thm:isdefeq-mono}
\leanok
One case per \texttt{PrimSpec} shape, seven in all. The spec's single lookup is pulled back along
\texttt{addConst\_eq\_of\_ne}, and everything else it asserts is monotone in the environment, so it
transfers by \texttt{VEnv.addConst\_le} or \texttt{VEnv.addDefEq\_le}. The bitwise shape is the
only one needing care, and only because its own quantification over larger environments has to be
composed with the new extension.
\end{proof}

\begin{theorem}[The single environment-extension theorem for the primitive invariant]
\label{thm:hasprimitives-addconst}
\lean{Lean4Lean.VEnv.HasPrimitives.addConstDefEq, Lean4Lean.primitives_eq, Lean4Lean.List.nodup_keys_unique, Lean4Lean.primSpecs_nodup, Lean4Lean.VEnv.HasPrimitives.addPrimitiveDefEq, Lean4Lean.mem_primSpecs_contains, Lean4Lean.VEnv.HasPrimitives.addConstDefEq_of_not_primitive, Lean4Lean.VEnv.HasPrimitives.addConst, Lean4Lean.VEnv.HasPrimitives.addDefEq}
\leanok
\uses{ind:primspec, def:kernel-primitives}
If \texttt{HasPrimitives env} holds and \texttt{env'} is \texttt{env} plus a constant, then
\texttt{env'.addDefEq df} again satisfies \texttt{HasPrimitives}, provided the caller discharges
the entry matching the new name. Because \texttt{primSpecs} is keyed by name without repetition
there is at most one such entry (\texttt{addPrimitiveDefEq}), and adding a name that is not a
reserved primitive discharges it vacuously.
\end{theorem}
\begin{proof}
\uses{thm:primspec-addconst}
\leanok
A dichotomy on whether the added name is the spec's key: the negative branch is
Theorem \ref{thm:primspec-addconst}, the positive branch is the caller's obligation.
\texttt{List.nodup\_keys\_unique} together with \texttt{primSpecs\_nodup}, proved by
\texttt{decide}, identifies the matching entry uniquely, and \texttt{primitives\_eq}
(\texttt{rfl}) is what lets a non-primitive name close the branch.
\end{proof}

\section{Bundled builders and conditional invariants}

\begin{definition}[A bundled translated-and-typed term]
\label{struct:trterm}
\lean{Lean4Lean.TrTerm, Lean4Lean.TrTy, Lean4Lean.TrTerm.of, Lean4Lean.TrTy.of, Lean4Lean.TrTerm.cast}
\leanok
\uses{ind:trexprs, def:hastype-istype}
\texttt{TrTerm env Us $\Delta$ src A} packages a model term \texttt{tgt}, a translation
\texttt{TrExprS env Us $\Delta$ src tgt}, and a typing of \texttt{tgt} at \texttt{A}; \texttt{TrTy}
is the analogue for types. The bundle is indexed by the \emph{source} expression, so the term the
checker built fixes the goal and the model term is an output. The module docstring states the
motivation: the two side conditions of the \texttt{TrExprS.app} rule are exactly a typing for the
function and a typing for the argument, which is the data such a bundle carries anyway, so the
primitive-constant proofs get their translations and typings by construction rather than by
threading them node by node.
\end{definition}

\begin{lemma}[Combinators for building translated terms]
\label{family:trterm-builders}
\lean{Lean4Lean.VLCtx.find?_vlam_self, Lean4Lean.VLCtx.find?_vlam_ne, Lean4Lean.VLCtx.find?_bvar_ne, Lean4Lean.TrExprS.of_nil, Lean4Lean.TrTerm.of_nil', Lean4Lean.TrTerm.natZero, Lean4Lean.TrTerm.boolLit, Lean4Lean.TrTerm.natSucc, Lean4Lean.TrTerm.wk, Lean4Lean.TrTerm.app, Lean4Lean.TrTerm.app', Lean4Lean.TrTerm.fvar, Lean4Lean.TrTerm.const0, Lean4Lean.TrTerm.natBinApp, Lean4Lean.TrTerm.boolBinApp, Lean4Lean.TrTerm.natUnApp, Lean4Lean.TrTy.mdata, Lean4Lean.TrTerm.mdata, Lean4Lean.TrTerm.lam, Lean4Lean.TrTy.forallE}
\leanok
\uses{struct:trterm, family:trexprs-literals}
The leaves are the three \texttt{find?} lemmas (a binder's own variable, a variable from an
earlier binder, a variable past an anonymous binder) and \texttt{of\_nil}, which transports a
translation established at the empty local context into any well-formed one provided its target is
closed. On top of them sit the application, lambda, pi and \texttt{mdata} combinators, specialised
to the non-dependent shapes the primitive probes use so that the side conditions are discharged by
\texttt{rfl}, and \texttt{wk}, which weakens a bundle past a freshly introduced binder.
\end{lemma}
\begin{proof}
\uses{thm:trexprs-weakfv}
\leanok
The \texttt{find?} lemmas unfold \texttt{VLCtx.find?} and \texttt{next}; \texttt{of\_nil} is
Theorem \ref{thm:trexprs-weakfv} along \texttt{FVLift.from\_nil} followed by
\texttt{ClosedN.liftN\_eq}; each combinator is then a record literal over the corresponding
\texttt{TrExprS} rule.
\end{proof}

\begin{theorem}[A primitive used as an operator enters as a bundle]
\label{thm:trexprs-ofconst}
\lean{Lean4Lean.TrExprS.ofConst, Lean4Lean.TrTerm.ofConst}
\leanok
\uses{ind:trexprs, struct:trterm, struct:orderedstrong}
\contrib{mixed} \srcloc{Lean4Lean/Verify/Typing/TrTerm.lean}{104}
\thesisref{no counterpart}
From a closed typing of \texttt{const c []} at \texttt{A} in the empty context over an
\texttt{OrderedStrong} environment, and given \texttt{A.instL [] = A} (discharged by \texttt{simp}
for the closed arrow types the primitives use), \texttt{const c []} translates to itself and has
type \texttt{A} in any level context and any local context.
\end{theorem}
\begin{proof}
\uses{fam:strong-inversions, thm:isdefeq-instl}
\leanok
\texttt{HasType.const\_inv} recovers the constant's declaration and universe-parameter count, which
is what the \texttt{const} rule of \texttt{TrExprS} needs; \texttt{instL} at the empty level list
and \texttt{weak0} then move the typing into the ambient context.
\end{proof}

\begin{remark}
The hypothesis of both members was changed on \texttt{iota} from \texttt{env.Ordered} to
\texttt{env.OrderedStrong}, forced by the new signature of \texttt{HasType.const\_inv}
(\srcloc{Lean4Lean/Theory/Typing/Strong.lean}{915}); the proof terms are otherwise identical to
master's. These two lines are the whole of the \texttt{iota} branch's footprint in
\srcloc{Lean4Lean/Verify/Typing/TrTerm.lean}{1}.
\end{remark}

\begin{definition}[Conditionally translatable expressions]
\label{def:conditionallytyped}
\lean{Lean4Lean.ConditionallyTyped, Lean4Lean.ConditionallyTyped.mk, Lean4Lean.ConditionallyTyped.mono, Lean4Lean.ConditionallyTyped.weakN_inv, Lean4Lean.ConditionallyTyped.fresh}
\leanok
\uses{ind:trexprs, def:ngen-reserves, def:vt-closed, def:vt-fvarsin}
\texttt{ConditionallyTyped ngen env Us $\Delta$ e} says that \texttt{e} is \texttt{Closed}, that
every free variable of \texttt{e} is reserved by the name generator, and that \emph{if} all its
free variables are in $\Delta$\texttt{.fvars} then \texttt{e} is translatable there. This is the
invariant the checker's state carries about terms whose context has been narrowed, so that a term
may survive a context it no longer fits. \texttt{mono} transports the invariant along generator
growth, \texttt{fresh} extends the context by a fresh declaration, and \texttt{weakN\_inv} removes
a named entry again.
\end{definition}
\axfoot{sorryAx in weakN\_inv only, via TrExprS.weakFV\_inv and TrProj.weak'\_inv}

\begin{definition}[Conditionally typed expressions]
\label{def:conditionallyhastype}
\lean{Lean4Lean.ConditionallyHasType, Lean4Lean.ConditionallyHasType.mk, Lean4Lean.ConditionallyHasType.typed, Lean4Lean.ConditionallyHasType.mono, Lean4Lean.ConditionallyHasType.weakN_inv, Lean4Lean.ConditionallyHasType.fresh}
\leanok
\uses{def:conditionallytyped, def:isfvarupset, thm:trexprs-weakfv-inv}
The same conditional discipline applied to a pair, a term and its type: both are closed and
reserved, and if the term's free variables are in the context then there is a full
\texttt{TrTyping} for the pair. \texttt{weakN\_inv} is the intricate member: it strengthens both
components with Theorem \ref{thm:trexprs-weakfv-inv}, realigns them with unique typing, pushes the
model typing down with \texttt{HasType.weakN\_iff}, and rebuilds the \texttt{FVarsBelow} clause
against the up-set relativised to the smaller context. It is sorry-dependent through that
strengthening.
\end{definition}
\axfoot{sorryAx in weakN\_inv, via TrExprS.weakFV\_inv, TrProj.weak'\_inv and VEnv.IsDefEqU.weakN\_iff}

\begin{definition}[Conditionally reduced expressions]
\label{def:conditionallywhnf}
\lean{Lean4Lean.ConditionallyWHNF, Lean4Lean.ConditionallyWHNF.mk, Lean4Lean.ConditionallyWHNF.mono, Lean4Lean.ConditionallyWHNF.weakN_inv, Lean4Lean.ConditionallyWHNF.fresh}
\leanok
\uses{def:conditionallyhastype, thm:isdefequ-weakn-iff}
The conditional discipline for a term together with its weak head normal form: both closed and
reserved, and, when the term's variables are in the context, a \texttt{FVarsBelow} and a shared
translation target for the two. \texttt{weakN\_inv} follows the same script as for
\texttt{ConditionallyHasType}, except that the fact pushed down at the end is an
\texttt{IsDefEqU.weakN\_iff} rather than a \texttt{HasType.weakN\_iff}; it inherits the same
admissions.
\end{definition}
\axfoot{sorryAx in weakN\_inv, via TrExprS.weakFV\_inv and VEnv.IsDefEqU.weakN\_iff (Theory/Typing/UniqueTyping.lean:172)}

\section{Contribution summary}

The \texttt{trproj} branch contributed the projection translation. Where
\srcloc{Lean4Lean/Verify/Typing/Expr.lean}{69} to 136 now stands, master held the single line
\texttt{def TrProj : ... := sorry}, and \srcloc{Lean4Lean/Verify/Typing/Lemmas.lean}{641} held seven
\texttt{sorry} lemmas about it, so the \texttt{proj} rule of \texttt{TrExprS}
(Definition \ref{ind:trexprs}) related a source projection to an arbitrary model term and the whole
translation relation was, in that case, vacuous. The branch replaced the stub by
\texttt{TrProjCtor} (Definition \ref{struct:trprojctor}) and its existential closure
\texttt{TrProj} (Definition \ref{def:trproj}): since \texttt{VExpr} has no projection node, the
target is required to be the recursor expansion applied to the major premise, which is the
thesis's \texttt{inv\_x} construction generalised from \texttt{Acc} to an arbitrary
single-constructor structure, built from the \texttt{Theory/Proj.lean} vocabulary of
Chapter \ref{chap:proj}. Two design decisions carry the section: the expansion's typing is a
\emph{field} of the relation rather than something to be derived, and the expansion data
(constructor name, level lists, parameters, field telescope) is hidden existentially so that the
consumers' statements do not change.

Five of master's seven stubs are now proved: Theorems \ref{thm:trproj-weak},
\ref{thm:trproj-defeqdfc}, \ref{thm:trproj-wf}, \ref{thm:trproj-instn} and
\ref{thm:trproj-instl}. Three of them, \ref{thm:trproj-weak}, \ref{thm:trproj-instn} and
\ref{thm:trproj-instl}, are a single \texttt{refine} over a record literal built on one observation
reused three times, that the pattern key, the field index and the minor arity mention the field
telescope only through its length; \ref{thm:trproj-defeqdfc} reuses the old record with
\texttt{\{H with \ldots\}} and \ref{thm:trproj-wf} is two lines. Theorem \ref{thm:trproj-wf} is
strictly stronger than master's stub, the hypothesis on the major premise being now a field; as a
result the projection case of Theorem \ref{thm:trexprs-wf} shrinks rather than grows. The new
parameters on \texttt{TrProj} create one obligation that master did not have,
Theorem \ref{thm:trproj-mono}, which is discharged using the \texttt{iota} branch's monotonicity
of the pattern registry. Two of the five proved forms also carry hypotheses master's stubs did not:
\texttt{Ordered env} on \ref{thm:trproj-weak} and, on \ref{thm:trproj-instn}, that plus a typing
for the substituted term; each is available at the single call site. Outside the projection section
the branch also added the three spine lemmas of Theorem \ref{thm:mkapplist-spine}, one of which is
used at \srcloc{Lean4Lean/Verify/TypeChecker/WHNF.lean}{50}, and the inversion
Theorem \ref{thm:trexprs-mkapplist-inv}, used five times at \texttt{WHNF.lean:87} to \texttt{:93},
both in the projection and $\iota$ proof. Downstream, the projection relation is what
\ref{thm:tr-env-proj-defeq} in Chapter \ref{chap:trenv} supplies the kernel-side reduction facts
for, and what \texttt{inferProj} in Chapter \ref{chap:typechecker} is meant to produce.

The \texttt{iota} branch contributed two unrelated things here. The first is the self-contained
block of Lemma \ref{family:lctx-empty-decl} and Theorem \ref{thm:mkforall-eq-fold}: lookups in a
chain of \texttt{mkLocalDecl} steps over the empty context, with the freshness side condition
deliberately dropped so that a six-step chain does not generate six obligations, five of its
eight declarations being consumed by the hand-built \texttt{Quot} axiom shapes in
\srcloc{Lean4Lean/Verify/Environment/Quot.lean}{1} and the rest serving them locally. The second is not a contribution of content
but a propagation: eight lemma signatures in the literal family
(Lemma \ref{family:trexprs-literals}) and two in Theorem \ref{thm:trexprs-ofconst} were
strengthened from \texttt{Ordered} to \texttt{OrderedStrong}, forced by the new signature of
\texttt{HasType.const\_inv} on that branch.

What remains open. Theorem \ref{thm:trproj-weak-inv} and Theorem \ref{thm:trproj-uniq} are still
\texttt{sorry}. The first is blocked on a master admission, \texttt{VEnv.IsDefEqU.weakN\_iff}, and
keeps Theorem \ref{thm:trexprs-weakfv-inv} and the three \texttt{weakN\_inv} lemmas of
Definitions \ref{def:conditionallytyped}, \ref{def:conditionallyhastype} and
\ref{def:conditionallywhnf} sorry-dependent, exactly as on master. The second is blocked on
facts about the relation itself, and keeps Theorem \ref{thm:trexprs-uniq},
Lemma \ref{family:trexpr-builders} and everything built on them sorry-dependent; its docstring
does not address how the per-field level list is pinned, which may mean the relation is
under-constrained rather than the proof merely unwritten. Beyond this chapter, nothing constructs
a \texttt{TrProjCtor} witness in general: \texttt{inferProj.WF\_struct} is \texttt{sorry}, so the
relation is inhabited only by the hand-built examples of Chapter \ref{chap:proj}.

\section{Review notes}

\begin{itemize}
\item \textbf{Medium.} \texttt{TrProj.uniq} is still \texttt{sorry}
(\srcloc{Lean4Lean/Verify/Typing/Lemmas.lean}{992}). Its docstring lists what a proof needs, unique
typing of the projection function, type-former injectivity for the structure's level list, its
parameters and the head names, and functionality of the $\iota$ registry for the parameter count
and the field count, but it does not say how the per-field level list \texttt{uss} is pinned. Two
\texttt{TrProjCtor} witnesses for the same projection may choose different \texttt{uss}, in which
case the two expansions are syntactically different recursor applications at different universe
instances and a definitional equality between them is not obviously available. Either the sketch
is incomplete or \texttt{TrProjCtor} is under-constrained. \ref{thm:trexprs-uniq} and the
\texttt{proj} builder of \ref{family:trexpr-builders} depend on this hole.
\item \textbf{Medium.} \texttt{TrProj.weak'\_inv} is still \texttt{sorry}
(\srcloc{Lean4Lean/Verify/Typing/Lemmas.lean}{745}). The docstring correctly traces it to the
admitted master lemma \texttt{IsDefEqU.weakN\_iff}
(\srcloc{Lean4Lean/Theory/Typing/UniqueTyping.lean}{172}), but the consequence is that
\ref{thm:trexprs-weakfv-inv} and all three \texttt{Conditionally} strengthening lemmas
(\ref{def:conditionallytyped}, \ref{def:conditionallyhastype}, \ref{def:conditionallywhnf}) remain
sorry-dependent, and those are a large part of the checker's specification.
\item \textbf{Medium.} \texttt{TrProjCtor.pat} (\srcloc{Lean4Lean/Verify/Typing/Expr.lean}{97})
asserts only that \emph{some} reduct is registered under the $\iota$ key; the reduct itself is not
pinned, so \ref{def:trproj} alone does not determine reduction behaviour, and the key, being a
sum, does not exclude a different motive/minor/index split adding up to the same number. The
docstring acknowledges this and defers to \ref{thm:tr-env-proj-defeq} for the kernel-side facts,
but the relation is weaker than it looks, and this is the stated obstacle to
\ref{thm:trproj-uniq}.
\item \textbf{Low.} \texttt{TrProj} (\srcloc{Lean4Lean/Verify/Typing/Expr.lean}{133})
existentially quantifies six pieces of data, two of which are pinned by other fields: the
parameter count is the length of the parameter list by \texttt{params\_length}, and the field
telescope is determined by the constructor name, the level list and the parameters by
\texttt{ctor}. Replacing them by definitions would shrink the relation and remove an obligation
from each of the three proofs that rebuild the record (\texttt{weak'}, \texttt{instN},
\texttt{instL}), where they reappear as \texttt{by simpa} and \texttt{by rw [hlen]}.
\item \textbf{Low.} The \texttt{iota} block installs a global \texttt{DecidableEq FVarId} instance
inside \texttt{namespace Lean.LocalContext}
(\srcloc{Lean4Lean/Verify/LocalContext.lean}{183}), under a section heading about the empty local
context. The instance has nothing to do with \texttt{LocalContext}, and given the project's own
\texttt{LawfulBEq FVarId} (\srcloc{Lean4Lean/Verify/Expr.lean}{12}) it could have been
\texttt{instDecidableEqOfLawfulBEq}, as is already done for \texttt{DefinitionSafety} at
\srcloc{Lean4Lean/Verify/Expr.lean}{51}. No conflict exists today, since Lean core has no such
instance, but the placement invites one.
\item \textbf{Low.} \texttt{TrExprS.mkAppList\_inv}
(\srcloc{Lean4Lean/Verify/Typing/Lemmas.lean}{2349}) overlaps with the pre-existing
\texttt{AppStack.build} (\srcloc{Lean4Lean/Verify/Typing/Lemmas.lean}{2344}): both invert a
translated application spine. The pointwise-list shape is what the WHNF proof wants, so the
duplication is defensible, but it leaves two spine-inversion idioms in one file.
\item \textbf{Low.} Of the three spine lemmas the \texttt{trproj} branch added to
\srcloc{Lean4Lean/Verify/Expr.lean}{745}, only \texttt{getAppArgsList\_mkAppList} has an explicit
consumer (\srcloc{Lean4Lean/Verify/TypeChecker/WHNF.lean}{50});
\texttt{getAppFn\_mkAppList} is reached only as a \texttt{@[simp]} lemma, and
\texttt{getAppArgsRevList\_mkAppList} (\srcloc{Lean4Lean/Verify/Expr.lean}{749}) has no consumer
outside its own file, existing only to prove the forward form.
\item \textbf{Low.} The \texttt{iota} branch strengthened eight master lemma signatures in
\srcloc{Lean4Lean/Verify/Typing/Lemmas.lean}{1856} and two in
\srcloc{Lean4Lean/Verify/Typing/TrTerm.lean}{104} from \texttt{env.Ordered} to
\texttt{env.OrderedStrong}, forced by the new signature of \texttt{HasType.const\_inv}. Since
\texttt{OrderedStrong} is obtained only through \texttt{VEnv.WF.orderedStrong}, which invokes the
admitted \texttt{VEnv.WF.patsStrong} (\srcloc{Lean4Lean/Theory/Typing/EnvLemmas.lean}{334}), every
downstream use of the literal lemmas (\srcloc{Lean4Lean/Verify/TypeChecker/Reduce.lean}{16},
\srcloc{Lean4Lean/Verify/TypeChecker/IsDefEq.lean}{406},
\srcloc{Lean4Lean/Verify/TypeChecker/InferType.lean}{420}) now depends on that admission.
\item \textbf{Low.} The scope paragraph of \texttt{TrProjCtor}
(\srcloc{Lean4Lean/Verify/Typing/Expr.lean}{84}) makes precise claims about the kernel's
\texttt{inferProj}, which structures it accepts and when its \texttt{Prop} gate fires, that are
prose only: nothing in the repository formalizes them, and both \texttt{inferProj.WF\_struct} and
\texttt{inferProj.WF} (\srcloc{Lean4Lean/Verify/TypeChecker/InferType.lean}{392},
\srcloc{Lean4Lean/Verify/TypeChecker/InferType.lean}{407}) are \texttt{sorry}. The completeness
boundary is therefore asserted, not verified.
\item \textbf{Low.} Within the \texttt{iota} block of
\srcloc{Lean4Lean/Verify/LocalContext.lean}{181}, \texttt{toList\_empty} and \texttt{WF.empty}
exist only to serve \texttt{find?\_empty} three lines later, and have no other consumer.
\end{itemize}
